\chapter{机械臂运动学}
\label{ch:arm-kin}

% ════════════════════════════════════════════════════════════════
%  机械臂建模、规划与控制部 · 第 1 章
%  经典理论无教材抽取源，按编写规范 §一.9 用 WebSearch/WebFetch 重建到复现级，
%  并完整写入正文（§一.8 自包含铁律），\cite 只标出处。
%  承上复用 P0（ch:rigid_body / ch:lie），对照 P7 四足（quad_kin）。
%  POE 主线取 lynch2017modern；DH/逆解取 spong2006robot。
%  实证血肉取 Robot arm_dm 六轴力矩臂（docs/Robot吸收_arm机械臂运控.md）。
% ════════════════════════════════════════════════════════════════

\begin{practice}[前置自测]
开始之前，先试答下列问题。若已能流畅作答，可只速读本章表格与速查卡；若卡住，正是本章要为你解决的。
\begin{enumerate}
  \item 给定六个关节角 $\mathbf{q}=(q_1,\dots,q_6)$，你能写出末端执行器在基座系下的位姿 $\mathbf{T}_{0e}(\mathbf{q})$ 吗？它和上一章（\cref{ch:rigid_body}）的齐次变换连乘是什么关系？
  \item 描述一条连杆相对前一连杆的位姿，最少需要几个标量参数？为什么 Denavit--Hartenberg 用四个而非六个？
  \item 旋量积（指数积，POE）正运动学公式 $\mathbf{T}=e^{[\mathcal{S}_1]\theta_1}\cdots e^{[\mathcal{S}_n]\theta_n}\mathbf{M}$ 里，每个 $\mathcal{S}_i$ 是什么？它和 \cref{ch:lie} 的 $\se(3)$ 元素、指数映射有何联系？
  \item 给定一个期望末端位姿，关节角的解唯一吗？什么时候无解、什么时候有多解（如“肘上 / 肘下”）？
  \item 末端速度 $\dot{\mathbf{x}}$ 与关节速度 $\dot{\mathbf{q}}$ 的线性关系矩阵叫雅可比 $\mathbf{J}$。当 $\det(\mathbf{J}\mathbf{J}^\top)\to0$ 时机器人发生了什么？为什么此时直接用 $\mathbf{J}^{-1}$ 会让关节速度“爆炸”？
  \item 一个 6 自由度臂去跟踪一个只有 3 维（位置）的目标，多出来的 3 个自由度去哪了？如何利用它们做避障而不破坏末端位置？
\end{enumerate}
\end{practice}

\section*{本章目标}
学完本章，你应当能够：
\begin{enumerate}
  \item \textbf{建立}串联机械臂的连杆--关节模型，用位姿链 $\mathbf{T}_{0n}=\prod_i\mathbf{T}_{i-1,i}$ 把“关节角 $\to$ 末端位姿”表达为齐次变换连乘（承 \cref{ch:rigid_body}）；
  \item \textbf{掌握}两套正运动学建模法——Denavit--Hartenberg 四参数法（标准 vs 改进）与旋量积（POE），并说清各自的优劣与适用场景；
  \item \textbf{推导}并理解 POE 公式如何由 \cref{ch:lie} 的 $\se(3)$ 指数映射、空间/物体旋量、伴随变换自然给出；
  \item \textbf{求解}逆运动学的解析（闭式）解：几何解耦法、球腕的 Pieper 准则与 6R 解耦三步、\emph{螺旋式} Paden--Kahan 子问题（在 POE 上几何反解）、以及通用偏置腕 6R 的 dialytic 消元法（解数尖锐上界 16），并讨论多解 / 无解 / 工作空间；
  \item \textbf{构造}几何雅可比（逐列=关节贡献）、解析雅可比、物体雅可比，理解静力学对偶 $\boldsymbol{\tau}=\mathbf{J}^\top\mathbf{F}$；
  \item \textbf{判别}奇异构型，用可操作度 $w=\sqrt{\det(\mathbf{J}\mathbf{J}^\top)}$ 与阻尼最小二乘（DLS）伪逆在奇异附近稳健求解；
  \item \textbf{运用}冗余与零空间投影 $\mathbf{N}=\mathbf{I}-\mathbf{J}^{+}\mathbf{J}$，用 resolved-rate 框架在执行主任务的同时做次要任务（避障 / 重构姿态）；
  \item \textbf{衔接}Pinocchio 等动力学库的 FK / Jacobian API，并把本章作为后续动力学（\cref{ch:arm-dyn}）、轨迹（\cref{ch:arm-traj}）、控制（\cref{ch:arm-ctrl}）章的几何地基。
\end{enumerate}

\subsection*{本章知识导航}
本章是一条“\emph{把关节空间与任务空间双向连接起来}”的主线：正运动学（关节 $\to$ 末端）、逆运动学（末端 $\to$ 关节）、速度层（雅可比）、以及雅可比退化（奇异）与富余（冗余）两个极端。结构如下：

\begin{center}
\small
\begin{tikzpicture}[
  node distance=4mm and 7mm, >=Stealth,
  blk/.style={draw, rounded corners, align=center, font=\small, inner sep=3pt, fill=main!6, draw=main!60!black},
  grp/.style={draw, rounded corners, align=center, font=\small, inner sep=3pt, fill=second!6, draw=second!60!black},
  ln/.style={->, thick, gray!60!black}]
\node[blk] (chain) {连杆位姿链\\ $\mathbf{T}_{0n}=\prod\mathbf{T}_{i-1,i}$};
\node[blk, right=of chain] (dh) {DH 四参数\\（坐标系法）};
\node[blk, right=of dh] (poe) {旋量积 POE\\（无坐标系法）};
\node[grp, below=8mm of chain] (ik) {逆运动学\\解析解 / Pieper};
\node[grp, right=of ik] (jac) {雅可比\\ $\dot{\mathbf{x}}=\mathbf{J}\dot{\mathbf{q}}$};
\node[grp, right=of jac] (sing) {奇异 + 可操作度\\ DLS 伪逆};
\node[blk, below=8mm of ik, fill=third!8, draw=third!60!black] (redun) {冗余与零空间\\ $\mathbf{N}=\mathbf{I}-\mathbf{J}^{+}\mathbf{J}$};
\node[blk, right=of redun, fill=third!8, draw=third!60!black] (prac) {实践指针\\ Pinocchio FK/Jacobian};
\draw[ln] (chain)--(dh); \draw[ln] (dh)--(poe);
\draw[ln] (poe.south) -- ++(0,-4.5mm) -| (jac.north);
\draw[ln] (chain.south) -- (ik.north);
\draw[ln] (jac)--(sing);
\draw[ln] (jac.south) -- (redun.north -| jac.south);
\draw[ln] (redun)--(prac);
\end{tikzpicture}
\end{center}

\noindent\textbf{推荐路径}：§1.1--1.3 是任何读者的主干（建立位姿链 + 两套正运动学建模法）；§1.4 逆运动学解析解是机械臂区别于一般刚体的招牌内容；§1.5--1.7（雅可比、奇异、冗余）是\emph{控制与规划所需}的核心工具，后续动力学（\cref{ch:arm-dyn}）、轨迹（\cref{ch:arm-traj}）、控制（\cref{ch:arm-ctrl}）章会反复引用——其中雅可比（\cref{sec:ak-jacobian}）是操作空间控制与力控的前置，DLS 伪逆（\cref{sec:ak-singularity}）是 P8 强化学习操作里 DiffIK 动作空间的数学背景（\cref{ch:rlmc-manip}）。§1.8 给出工程落地指针，指向本部实践章 \cref{ch:arm-prac}。

\subsection*{前置知识桥接}
\cref{ch:rigid_body} 已给出 $\SE(3)$ 位姿 $\mathbf{T}$、齐次变换的连乘与求逆（\cref{def:se3-rep,prop:T-inv}）、反对称算子 $(\cdot)^\wedge$ 与恒等式（\cref{sec:rb-vector}）、旋转/位姿运动学（\cref{eq:poisson}）。\cref{ch:lie} 进一步给出 $\so(3),\se(3)$（\cref{def:so3-se3}）、$\SE(3)$ 指数映射（\cref{thm:se3-exp}）、伴随 $\mathrm{Ad}(\mathbf{T})$（\cref{eq:Ad}）、以及位姿运动学的李代数形式 $\boldsymbol{\varpi}=\boldsymbol{\mathcal{J}}\dot{\boldsymbol{\xi}}$（\cref{thm:se3-kinematics}）。本章把这些一般刚体工具\emph{专门化}到“多关节串联链”这一结构上：FK 是齐次变换连乘的具体落实，POE 是 $\SE(3)$ 指数映射的连乘，雅可比是位姿运动学的逐关节展开。换言之，从“一个刚体如何动”进到“\emph{一串铰接刚体如何把关节运动传递到末端}”。

\subsection*{与 P7 四足部的对照}
\cref{ch:quad_kin}（四足运动学）把\emph{单腿}当作一条串联链处理，但因腿短、构型规整，它用\emph{解耦闭式}正逆运动学（\cref{sec:qk-fk,sec:qk-ik}）而非通用 DH；它的雅可比（\cref{sec:qk-jacobian}）与零空间理论（\cref{sec:qk-nullspace}）则与本章\emph{同源}。本章与之互为印证：\textbf{本章给定基串联臂的通用框架（DH/POE/解析 IK），P7 给规整短链的捷径}；零空间投影、MP 伪逆、阻尼最小二乘等数学内核在 \cref{sec:qk-nullspace} 已有完整证明，本章\emph{直接复用不重证}，只补“机械臂视角”（降维造冗余、resolved-rate、可操作度椭球）。

\begin{table}[H]\centering\small
\caption{本章阅读方式建议}
\begin{tabular}{lll}
\toprule
阅读方式 & 时间 & 适合谁 \\
\midrule
精读（含全部推导与练习） & 4--6 小时 & 首次系统学习机械臂、要做规控推导 \\
速读（跳过冗长推导细节） & 1.5 小时 & 有运动学基础、想对齐本书记号与 POE 主线 \\
速查（只看小结的符号表 + 速查表） & 15 分钟 & 写代码时回来查公式 \\
\bottomrule
\end{tabular}
\end{table}

\begin{note}[本章记号与约定]\label{note:ak-notation}
关节变量 $\mathbf{q}=(q_1,\dots,q_n)\in\mathbb{R}^n$（转动关节 $q_i=\theta_i$ 取角度、移动关节 $q_i=d_i$ 取位移）；连杆坐标系 $\{i\}$，第 $i$ 关节连接连杆 $i-1$ 与 $i$。位姿 $\mathbf{T}\in\SE(3)$（\cref{def:se3-rep}），$\mathbf{T}_{i-1,i}$ 表“$\{i\}$ 在 $\{i-1\}$ 中的位姿”。末端执行器（end-effector，EE）系记 $\{e\}$ 或 $\{n\}$，基座系 $\{0\}$。旋量 $\mathcal{S}=[\boldsymbol{\omega};\,\mathbf{v}]\in\mathbb{R}^6$（本书 $\se(3)$ 排序为\textbf{平移在前}，但旋量记号沿用 Modern Robotics 习惯把角速度部分 $\boldsymbol{\omega}$ 写在前、记 $\mathcal{S}=(\boldsymbol{\omega},\mathbf{v})$，二者只是分量排列、几何无别；用到 $\se(3)$ 向量与本书 $\boldsymbol{\xi}=[\boldsymbol{\rho};\boldsymbol{\phi}]$ 对接处显式说明）。$[\mathcal{S}]\in\se(3)$ 为 $\mathcal{S}$ 的 $4\times4$ 矩阵形式。雅可比 $\mathbf{J}\in\mathbb{R}^{6\times n}$（任务为全位姿）或 $\mathbb{R}^{m\times n}$（任务维 $m$）。粗体小写为向量、粗体大写为矩阵；力矩 $\boldsymbol{\tau}$、末端广义力（wrench）$\mathbf{F}=[\mathbf{m};\mathbf{f}]$（力矩在前、力在后，与旋量配对）。
\end{note}

% ════════════════════════════════════════════════════════════════
\section{连杆、关节与位姿链}
\label{sec:ak-chain}

\paragraph{动机：从“一个刚体”到“一串铰接刚体”。}
\cref{ch:rigid_body} 教会我们描述\emph{一个}刚体的位姿 $\mathbf{T}\in\SE(3)$ 及其变换。但机械臂是\emph{多个}刚体（连杆）用关节（铰）串起来的：每个关节给出一个自由度，整条链的末端位姿由所有关节角共同决定。运动学的第一任务，就是把这个“关节角 $\to$ 末端位姿”的映射\emph{写下来}。

\paragraph{反面：把整臂当一个黑箱会失去结构。}
我们当然可以把末端位姿看成关节角的某个抽象函数 $\mathbf{T}_{0e}=f(\mathbf{q})$，但若不利用“链式结构”，这个 $f$ 既写不出闭式、也无法求导。串联机械臂的关键结构是：\textbf{相邻连杆之间的相对位姿只依赖一个关节变量}，整链位姿是这些局部位姿的\emph{连乘}。抓住这一点，FK、雅可比、IK 全都顺理成章。

\begin{insight}[同一个 $4\times4$，四层读法：从代数对象到物体描述]
位姿链的砖块是相邻位姿 $\mathbf{T}_{i-1,i}\in\SE(3)$（\cref{def:se3-rep}）。为什么单单一个 $4\times4$ 矩阵就能撑起整套运动学？Paul 在最早系统使用齐次变换的著作里，给出一个层层加深的读法\cite{paul1981robot}——同一个矩阵可\emph{依次}读作四样东西，理解每深一层，后文（位姿链、DH、POE、雅可比）就多一分透彻：
\begin{enumerate}
  \item \textbf{纯代数对象}：它首先只是一个作用在齐次坐标 $[x,y,z,1]^\top$ 上的 $4\times4$ 矩阵，满足结合律、可求逆（\cref{prop:T-inv}）。这一层只问“怎么算”，不问“是什么”——但正是这层代数让“连乘”有意义。
  \item \textbf{坐标变换算子}：把一个点在内层系 $\{i\}$ 下的坐标，左乘 $\mathbf{T}_{i-1,i}$ 即得它在外层系 $\{i-1\}$ 下的坐标。此时矩阵是个\emph{动作}——“换一套坐标来描述同一个点”，这正是 \cref{thm:ak-poschain} 逐级把点从末端系映回基座系的微观机理。
  \item \textbf{一个坐标系的表示}：它的前三列是 $\{i\}$ 三根坐标轴在 $\{i-1\}$ 中的单位方向、第四列是 $\{i\}$ 原点的位置\cite{paul1981robot}。于是矩阵\emph{本身}就\emph{是}一个被摆放好的坐标系——这一层让 DH（\cref{sec:ak-dh}，给每根连杆“摆一个系”）和 POE 的零位 $\mathbf{M}$（\cref{sec:ak-poe}）都成了“摆系”的具体手艺。
  \item \textbf{一个物体的描述}：把坐标系\emph{固连}在连杆（刚体）上，矩阵便描述了\emph{那段连杆}的位置与姿态；且“$A$ 相对 $B$”可求逆得“$B$ 相对 $A$”\cite{paul1981robot}——这是普通的相对位移向量\emph{办不到}的。整条机械臂，正是 $n$ 个这样的物体描述首尾相扣。
\end{enumerate}
四层是\emph{同一矩阵}的四副面孔，不是四样东西：算它时把它当代数(1)，连乘时当算子(2)，建模放系时当坐标系(3)，谈连杆位姿时当物体(4)。本章后文会反复在这四层间切换——读到费解处，不妨自问“此刻我把这个 $\mathbf{T}$ 当哪一层在用？”
\end{insight}

\begin{definition}[串联机械臂与关节变量]\label{def:ak-serial}
\textbf{串联机械臂}（serial manipulator）是 $n+1$ 个刚体连杆 $L_0,L_1,\dots,L_n$ 经 $n$ 个单自由度关节依次连接的开链（open chain），$L_0$ 为固定基座、$L_n$ 末端装末端执行器。第 $i$ 个关节连接 $L_{i-1}$ 与 $L_i$，其\textbf{关节变量} $q_i$ 为：转动关节（revolute）的转角 $\theta_i$，或移动关节（prismatic）的位移 $d_i$。把每根连杆固连一个坐标系 $\{i\}$（$\{0\}$ 为基座系），则\textbf{相邻连杆相对位姿} $\mathbf{T}_{i-1,i}(q_i)$ 仅依赖 $q_i$ 一个变量。
\end{definition}

\begin{theorem}[正运动学的位姿链]\label{thm:ak-poschain}
$n$ 自由度串联机械臂的末端执行器在基座系下的位姿为相邻连杆位姿的连乘：
\begin{equation}
  \boxed{\mathbf{T}_{0n}(\mathbf{q})=\mathbf{T}_{01}(q_1)\,\mathbf{T}_{12}(q_2)\cdots\mathbf{T}_{n-1,n}(q_n)=\prod_{i=1}^{n}\mathbf{T}_{i-1,i}(q_i).}
  \label{eq:ak-poschain}
\end{equation}
\end{theorem}
\begin{proof}
由 $\SE(3)$ 变换的复合律（\cref{def:se3-rep} 后的连乘性质）：若 $\{i\}$ 在 $\{i-1\}$ 中位姿为 $\mathbf{T}_{i-1,i}$、$\{i+1\}$ 在 $\{i\}$ 中为 $\mathbf{T}_{i,i+1}$，则 $\{i+1\}$ 在 $\{i-1\}$ 中为 $\mathbf{T}_{i-1,i}\mathbf{T}_{i,i+1}$（齐次坐标下即矩阵乘法，逐级把点从内层系映回外层系）。归纳应用到全链 $\{0\}\to\{1\}\to\cdots\to\{n\}$ 即得 \cref{eq:ak-poschain}。若末端执行器系 $\{e\}$ 相对 $\{n\}$ 还有一个固定安装位姿 $\mathbf{T}_{ne}$（常量），则 $\mathbf{T}_{0e}=\mathbf{T}_{0n}\mathbf{T}_{ne}$。
\end{proof}

\begin{insight}[正运动学是“确定性可解”的，逆才是难题]
\cref{eq:ak-poschain} 揭示一个反差：给定 $\mathbf{q}$ 算 $\mathbf{T}_{0n}$（\textbf{正}运动学）永远是\emph{矩阵连乘}，确定、唯一、$O(n)$ 即可算完——再复杂的臂也不过多乘几个 $4\times4$。难的是反过来：给定 $\mathbf{T}_{0n}$ 求 $\mathbf{q}$（\textbf{逆}运动学），这是解一组非线性三角方程，可能无解、多解（\cref{sec:ak-ik-analytic}）。这个“正易逆难”的不对称，是机械臂运动学全部技术张力的来源。
\end{insight}

\begin{figure}[ht]
\centering
\begin{tikzpicture}[scale=1.0,>=Stealth,
  lk/.style={line width=2.2pt, main!70!black, cap=round},
  jt/.style={circle, draw=second!70!black, fill=second!20, inner sep=1.6pt},
  ax/.style={->, thick, third!70!black}]
% base
\draw[fill=gray!30, draw=gray!60] (-0.4,-0.25) rectangle (0.4,0); \node[below] at (0,-0.25) {\scriptsize 基座 $L_0$, 系 $\{0\}$};
% link chain
\node[jt] (j1) at (0,0) {}; \node[left] at (-0.15,0.05) {\scriptsize 关节 1 ($q_1$)};
\draw[lk] (j1) -- (1.4,0.9); \node[jt] (j2) at (1.4,0.9) {}; \node[above left=-1pt] at (j2) {\scriptsize 关节 2 ($q_2$)};
\draw[lk] (j2) -- (2.9,0.6); \node[jt] (j3) at (2.9,0.6) {}; \node[above] at (j3) {\scriptsize 关节 3};
\draw[lk] (j3) -- (4.0,1.3); \node[jt] (je) at (4.0,1.3) {};
\node[right] at (4.1,1.35) {\scriptsize 末端 $\{e\}$};
% frame triads (schematic)
\draw[ax] (j1) -- ++(0.5,0); \draw[ax] (j1) -- ++(0,0.5);
\draw[ax] (je) -- ++(0.45,0.15); \draw[ax] (je) -- ++(-0.15,0.45);
% link transforms labels
\node[main!60!black] at (0.7,0.65) {\scriptsize $\mathbf{T}_{01}$};
\node[main!60!black] at (2.2,1.0) {\scriptsize $\mathbf{T}_{12}$};
\node[main!60!black] at (3.55,0.75) {\scriptsize $\mathbf{T}_{23}$};
\end{tikzpicture}
\caption{串联机械臂的连杆--关节链与位姿链。每个关节变量 $q_i$ 决定一个相邻位姿 $\mathbf{T}_{i-1,i}(q_i)$，末端位姿是它们的连乘 \cref{eq:ak-poschain}。}
\label{fig:ak-chain}
\end{figure}

\paragraph{案例：六轴力矩臂的链结构。}
本部实践章（\cref{ch:arm-prac}）所用的 \texttt{arm\_dm} 六轴力矩臂是一条典型 6R 链：腰 $J_1$（绕 $z$）+ 肩 $J_2$、肘 $J_3$（绕 $y$）构成定位用的\emph{臂部}，腕 $J_4,J_5,J_6$（绕 $x,y,x$）构成定姿用的\emph{球腕}。这种“前三轴定位、后三轴定姿、腕三轴交于一点”的结构不是偶然——它正是为了让逆运动学\emph{可解耦出闭式}（\cref{sec:ak-ik-analytic} 的 Pieper 准则）。该臂\textbf{重肩肘、轻腕}的异质惯量特性（腕惯量仅约 $2\times10^{-4}$）虽属动力学范畴（\cref{ch:arm-dyn}），但其几何根源——腕是末端的小法兰——已在链结构中埋下。

\begin{practice}
\begin{enumerate}
  \item 写出一个 2R 平面臂（连杆长 $a_1,a_2$，关节角 $\theta_1,\theta_2$）的 $\mathbf{T}_{01},\mathbf{T}_{12}$，并连乘得 $\mathbf{T}_{02}$，验证末端位置 $(a_1\cos\theta_1+a_2\cos(\theta_1+\theta_2),\,a_1\sin\theta_1+a_2\sin(\theta_1+\theta_2))$。
  \item 解释为什么 \cref{eq:ak-poschain} 的连乘顺序\emph{不能}颠倒（提示：左乘 vs 右乘对应在哪个系里施加变换）。
  \item 一台 SCARA 臂有 RRPR 四关节，其中第三关节为移动关节。说明此时 $q_3=d_3$、$\mathbf{T}_{23}(d_3)$ 是纯平移，正运动学公式形式不变。
\end{enumerate}
\end{practice}

% ════════════════════════════════════════════════════════════════
\section{正运动学：Denavit--Hartenberg 参数法}
\label{sec:ak-dh}

\paragraph{动机：用最少参数描述一条连杆。}
\cref{eq:ak-poschain} 把 FK 化为连乘，但还剩一个问题：$\mathbf{T}_{i-1,i}(q_i)$ 怎么写？一般 $\SE(3)$ 位姿要 6 个参数，但相邻连杆系并非任意放置——若\emph{巧妙地}规定每个连杆系的放法，就能用更少参数描述。1955 年 Denavit 与 Hartenberg 给出了这套规约\cite{spong2006robot}：每条连杆只需 \textbf{4 个参数}（而非 6 个），代价是放坐标系要守规则。

\paragraph{核心思想：把相对位姿压缩成两转两移。}
DH 的关键洞察是：两条相邻关节轴（两条空间直线）之间的相对几何，由它们的\emph{公垂线}唯一确定（除非平行）。沿“前一关节轴 $z_{i-1}$ 转 + 移、再沿公垂线 $x_i$ 转 + 移”这\emph{四步}即可从 $\{i-1\}$ 到 $\{i\}$。这把一个 6 参数位姿压缩到 4 参数——少掉的 2 个自由度正是被“坐标系必须按规则放”消化掉了。

\begin{definition}[Denavit--Hartenberg 四参数]\label{def:ak-dh}
按标准（distal）DH 规约放置连杆系后，从 $\{i-1\}$ 到 $\{i\}$ 由四个参数描述\cite{spong2006robot}：
\begin{itemize}
  \item $\theta_i$（\textbf{关节角}）：绕 $z_{i-1}$ 从 $x_{i-1}$ 转到 $x_i$ 的角；
  \item $d_i$（\textbf{连杆偏距}）：沿 $z_{i-1}$ 从 $x_{i-1}$ 到 $x_i$ 的有向距离；
  \item $a_i$（\textbf{连杆长度}）：沿 $x_i$ 的公垂线长（$z_{i-1}$ 与 $z_i$ 之间）；
  \item $\alpha_i$（\textbf{连杆扭角}）：绕 $x_i$ 从 $z_{i-1}$ 转到 $z_i$ 的角。
\end{itemize}
对转动关节 $\theta_i$ 是变量、其余为常量；对移动关节 $d_i$ 是变量、其余常量。
\end{definition}

\begin{theorem}[标准 DH 变换矩阵]\label{thm:ak-dh-matrix}
按 \cref{def:ak-dh} 的四参数，相邻连杆位姿为“先绕 $z_{i-1}$ 转 $\theta_i$ 移 $d_i$，再绕 $x_i$ 转 $\alpha_i$ 移 $a_i$”的复合\cite{spong2006robot}：
\begin{equation}
  \mathbf{T}_{i-1,i}=\mathrm{Rot}_z(\theta_i)\,\mathrm{Trans}_z(d_i)\,\mathrm{Trans}_x(a_i)\,\mathrm{Rot}_x(\alpha_i),
  \label{eq:ak-dh-compose}
\end{equation}
展开为
\begin{equation}
  \boxed{\mathbf{T}_{i-1,i}=
  \begin{bmatrix}
    \cos\theta_i & -\sin\theta_i\cos\alpha_i & \sin\theta_i\sin\alpha_i & a_i\cos\theta_i\\
    \sin\theta_i & \cos\theta_i\cos\alpha_i & -\cos\theta_i\sin\alpha_i & a_i\sin\theta_i\\
    0 & \sin\alpha_i & \cos\alpha_i & d_i\\
    0 & 0 & 0 & 1
  \end{bmatrix}.}
  \label{eq:ak-dh-matrix}
\end{equation}
\end{theorem}
\begin{proof}
逐个写出四个基本变换的 $4\times4$ 矩阵：$\mathrm{Rot}_z(\theta_i)$ 是绕 $z$ 转 $\theta_i$、$\mathrm{Trans}_z(d_i)$ 沿 $z$ 移 $d_i$、$\mathrm{Trans}_x(a_i)$ 沿 $x$ 移 $a_i$、$\mathrm{Rot}_x(\alpha_i)$ 绕 $x$ 转 $\alpha_i$（旋转部分见 \cref{ch:rigid_body} 的基本旋转矩阵）。按 \cref{eq:ak-dh-compose} 顺序相乘——注意均为\emph{右乘}（在当前动系里依次施加）——合并旋转块得 $\mathbf{R}=\mathrm{Rot}_z(\theta_i)\mathrm{Rot}_x(\alpha_i)$、平移块由前两步的 $z$ 平移 $d_i$ 与第三步的 $x$ 平移 $a_i$（经 $\mathrm{Rot}_z(\theta_i)$ 旋到 $(a_i\cos\theta_i,a_i\sin\theta_i,0)$）叠加，即 \cref{eq:ak-dh-matrix}。
\end{proof}

\paragraph{坐标系放置规则（标准 DH）。}
要让 \cref{eq:ak-dh-matrix} 成立，连杆系须按下述规则放（这是 DH 用得对不对的关键）\cite{spong2006robot}：
\begin{enumerate}
  \item \textbf{$z_i$ 沿第 $i+1$ 关节轴}（转动关节的转轴 / 移动关节的移动方向）；
  \item \textbf{$x_i$ 沿 $z_{i-1}$ 与 $z_i$ 的公垂线}，方向从 $z_{i-1}$ 指向 $z_i$；若两轴相交则取 $x_i\parallel z_{i-1}\times z_i$；若平行则公垂线不唯一，$x_i$ 可任选（通常令 $d_i=0$）；
  \item \textbf{原点 $O_i$} 取在 $x_i$ 与 $z_i$ 的交点；
  \item \textbf{$y_i$} 由右手系 $y_i=z_i\times x_i$ 补全。
\end{enumerate}
这样放，\cref{def:ak-dh} 的四个量恰好刻画 $\{i-1\}\to\{i\}$，且 $\mathbf{T}_{0n}=\prod_i\mathbf{T}_{i-1,i}$。整臂只需列一张 \textbf{DH 参数表}（每行四个数 + 关节类型）即可完全确定 FK。

\begin{note}[标准 DH vs 改进 DH（Craig 约定）]\label{note:ak-dh-modified}
存在两套 DH。\textbf{标准（distal）DH}把系 $\{i\}$ 放在连杆 $i$ 的\emph{远端}（关节 $i+1$ 处），矩阵即 \cref{eq:ak-dh-matrix}。\textbf{改进（proximal / Craig）DH}把系 $\{i\}$ 放在连杆 $i$ 的\emph{近端}（关节 $i$ 处），变换顺序改为“先绕 $x_{i-1}$ 转 $\alpha_{i-1}$ 移 $a_{i-1}$，再绕 $z_i$ 转 $\theta_i$ 移 $d_i$”\cite{spong2006robot}：
\[
  \mathbf{T}_{i-1,i}=
  \begin{bmatrix}
    \cos\theta_i & -\sin\theta_i & 0 & a_{i-1}\\
    \sin\theta_i\cos\alpha_{i-1} & \cos\theta_i\cos\alpha_{i-1} & -\sin\alpha_{i-1} & -d_i\sin\alpha_{i-1}\\
    \sin\theta_i\sin\alpha_{i-1} & \cos\theta_i\sin\alpha_{i-1} & \cos\alpha_{i-1} & d_i\cos\alpha_{i-1}\\
    0 & 0 & 0 & 1
  \end{bmatrix}.
\]
注意此时扭角/连杆长带的是\emph{前一连杆}的下标 $i-1$。两套\emph{描述同一臂}、FK 结果相同，但参数表的数值与下标对应\textbf{不同}——\textbf{混用是新手最常见的 DH 错误}。改进 DH 对分支链（多个子链共享祖先连杆）更自然，故 ROS/URDF 生态、Craig 教材、多数现代软件用改进 DH；Spong 等经典教材用标准 DH。本书用到具体参数表时显式声明用哪一套。
\end{note}

\begin{pitfall}[概念误区：DH 参数“唯一”]
新手以为一台臂的 DH 表是唯一的。实际上：(1) 标准 vs 改进两套就给出两张不同的表；(2) 即便同一套，平行轴 / 相交轴处公垂线不唯一、零位的选择有自由度，故 DH 表\emph{不唯一}。这带来工程陷阱：从厂商拿到 DH 表，必须\textbf{确认它是哪套约定、零位如何定义}，否则 FK 会系统性偏差。更稳妥的现代做法是用 URDF（每关节给出相对父连杆的完整 $\SE(3)$ 安装位姿 + 轴向），它\emph{回避}了 DH 的约定歧义——这也是下一节 POE 受青睐的部分原因。
\end{pitfall}

\paragraph{案例：URDF 而非 DH。}
值得注意的是，\cref{ch:arm-prac} 的 \texttt{arm\_dm} 臂\emph{不}用 DH 建模，而是用 SolidWorks 导出的 URDF——每个关节直接给出相对父连杆的 $\SE(3)$ 安装位姿与轴向单位向量。这正是工程现实：现代机器人软件栈（ROS、MoveIt、Pinocchio）以 URDF 为通用模型格式，DH 更多作为\emph{解析逆运动学推导}和\emph{教学}的工具存在。URDF 每关节的 $\langle\text{origin}\rangle$（父子系间固定变换）+ $\langle\text{axis}\rangle$（关节轴）恰好是下一节 POE 旋量所需的全部信息——这暗示 POE 比 DH 更贴合 URDF。

\begin{practice}
\begin{enumerate}
  \item 列出一个 3R 拟人臂（PUMA 前三轴：腰转 + 肩 + 肘）的标准 DH 参数表（自定连杆长），写出三个 $\mathbf{T}_{i-1,i}$ 并连乘。
  \item 把上题改用改进 DH 重列一遍，验证 $\mathbf{T}_{03}$ 与标准 DH 结果一致（虽中间矩阵不同）。
  \item 解释 \cref{eq:ak-dh-matrix} 里为何第 3 行第 1、2 列恒为 $(0,\sin\alpha_i)$、$(0,\cos\alpha_i)$——这反映了 DH 的哪条放置规则？
\end{enumerate}
\end{practice}

% ════════════════════════════════════════════════════════════════
\section{正运动学：旋量积（指数积，POE）}
\label{sec:ak-poe}

\paragraph{动机：免去逐连杆放坐标系的繁琐。}
DH 虽省参数，却要给\emph{每条}连杆精心放一个坐标系、还要分清两套约定——这既繁琐又易错。旋量积（product of exponentials，POE）给出另一条路\cite{lynch2017modern}：\textbf{只需两个坐标系}（基座系 $\{0\}$ 与末端零位系 $\{n\}$），加上每个关节的\emph{旋量}（screw axis）。它把 FK 写成 \cref{ch:lie} 的 $\SE(3)$ 指数映射的连乘，几何意义透明、统一处理转动 / 移动关节。

\paragraph{核心思想：每个关节是一次螺旋运动。}
回忆 \cref{ch:lie}：$\se(3)$ 元素经指数映射 $\exp([\mathcal{S}]\theta)$ 给出一个 $\SE(3)$ 位姿，几何上是绕某条空间轴的\emph{螺旋（旋拧）运动}（Chasles 定理）。转动关节 $i$ 转 $\theta_i$，正是末端绕该关节轴做一次纯旋转——一个螺旋运动 $\exp([\mathcal{S}_i]\theta_i)$。把全部关节的螺旋运动\emph{依次叠加}，再乘上零位位姿，就得到末端位姿。这就是 POE。

\begin{definition}[关节旋量（螺旋轴）]\label{def:ak-screw}
关节 $i$ 在基座系 $\{0\}$ 下的\textbf{空间旋量}（spatial screw axis）$\mathcal{S}_i=[\boldsymbol{\omega}_i;\,\mathbf{v}_i]\in\mathbb{R}^6$（这里按 Modern Robotics 习惯角速度分量在前）刻画该关节在\emph{零位}时引起的单位螺旋运动\cite{lynch2017modern}：
\begin{itemize}
  \item 转动关节：$\boldsymbol{\omega}_i$ 为关节轴单位向量，$\mathbf{v}_i=-\boldsymbol{\omega}_i\times\mathbf{q}_i$，其中 $\mathbf{q}_i$ 是轴上任一点（零位时在 $\{0\}$ 下的坐标）；
  \item 移动关节：$\boldsymbol{\omega}_i=\mathbf{0}$，$\mathbf{v}_i$ 为移动方向单位向量。
\end{itemize}
其矩阵形式 $[\mathcal{S}_i]=\big[\begin{smallmatrix}\boldsymbol{\omega}_i^\wedge & \mathbf{v}_i\\ \mathbf{0}^\top & 0\end{smallmatrix}\big]\in\se(3)$（\cref{def:so3-se3}）。
\end{definition}

\begin{note}[旋量记号与本书 $\se(3)$ 排序的对接]\label{note:ak-screw-order}
本书 $\se(3)$ 向量按\textbf{平移在前}排 $\boldsymbol{\xi}=[\boldsymbol{\rho};\boldsymbol{\phi}]$（\cref{note:ak-notation}、\cref{ch:lie}），而旋量文献（Lynch--Park）惯把\textbf{角速度在前}写 $\mathcal{S}=[\boldsymbol{\omega};\mathbf{v}]$。二者只是\emph{分量排列}不同，所指的 $\se(3)$ 元素、矩阵 $[\mathcal{S}]$、指数映射 $\exp([\mathcal{S}]\theta)$ 完全一致——用一个置换矩阵 $\mathbf{P}=\big[\begin{smallmatrix}\mathbf{0}&\mathbf{I}\\\mathbf{I}&\mathbf{0}\end{smallmatrix}\big]$ 即可互换（$\boldsymbol{\xi}=\mathbf{P}\mathcal{S}$）。本章为贴合旋量文献的几何叙述用 $[\boldsymbol{\omega};\mathbf{v}]$，凡与本书伴随 $\mathrm{Ad}$、雅可比对接处显式提醒排序。\textbf{切忌}把两套排序的分量直接代入对方公式。
\end{note}

\begin{theorem}[空间形式的旋量积 FK]\label{thm:ak-poe-space}
设末端在\textbf{零位}（所有 $\theta_i=0$）时位姿为 $\mathbf{M}\in\SE(3)$，关节空间旋量为 $\mathcal{S}_1,\dots,\mathcal{S}_n$（\cref{def:ak-screw}），则正运动学为\cite{lynch2017modern}：
\begin{equation}
  \boxed{\mathbf{T}_{0n}(\boldsymbol{\theta})=e^{[\mathcal{S}_1]\theta_1}\,e^{[\mathcal{S}_2]\theta_2}\cdots e^{[\mathcal{S}_n]\theta_n}\,\mathbf{M}.}
  \label{eq:ak-poe-space}
\end{equation}
其中 $e^{[\mathcal{S}_i]\theta_i}=\exp([\mathcal{S}_i]\theta_i)\in\SE(3)$ 是 \cref{thm:se3-exp} 的指数映射。
\end{theorem}
\begin{proof}
对 $n$ 用归纳法，体现 POE 的几何直觉\cite{lynch2017modern}。先看只动最后一个关节 $n$：其余关节锁在零位，末端绕关节 $n$ 的固定空间轴 $\mathcal{S}_n$ 做螺旋运动，由 \cref{ch:lie} 螺旋运动的指数公式，末端位姿 $=e^{[\mathcal{S}_n]\theta_n}\mathbf{M}$（在 $\mathbf{M}$ 上\emph{左乘}空间系下的螺旋，因 $\mathcal{S}_n$ 表达在 $\{0\}$ 中）。再加动关节 $n-1$：它的螺旋 $\mathcal{S}_{n-1}$ 也表达在 $\{0\}$ 中（零位定义），且关节 $n-1$ 在\emph{运动学上位于}关节 $n$ 与基座之间——它带动“关节 $n$ 及末端”这整个子刚体一起绕 $\mathcal{S}_{n-1}$ 螺旋，故再左乘 $e^{[\mathcal{S}_{n-1}]\theta_{n-1}}$。依次往基座方向递推，得 \cref{eq:ak-poe-space}。\emph{关键}：每个 $\mathcal{S}_i$ 是\emph{零位时}该关节轴的旋量，不随其它关节运动改变其在公式中的形式——这正是 POE 只需“零位 + 旋量”的原因。
\end{proof}

\begin{theorem}[物体形式的旋量积 FK]\label{thm:ak-poe-body}
等价地，可把关节旋量表达在\emph{末端}（物体）系 $\{n\}$ 中，记 $\mathcal{B}_i$，则\cite{lynch2017modern}：
\begin{equation}
  \mathbf{T}_{0n}(\boldsymbol{\theta})=\mathbf{M}\,e^{[\mathcal{B}_1]\theta_1}\,e^{[\mathcal{B}_2]\theta_2}\cdots e^{[\mathcal{B}_n]\theta_n},
  \label{eq:ak-poe-body}
\end{equation}
其中物体旋量由空间旋量经伴随变换得到：$\mathcal{B}_i=\mathrm{Ad}_{\mathbf{M}^{-1}}\mathcal{S}_i$（伴随 $\mathrm{Ad}$ 见 \cref{eq:Ad}）。
\end{theorem}
\begin{proof}
由 \cref{ch:lie} 伴随的核心恒等式 $\mathbf{M}\,e^{[\mathcal{A}]\theta}\,\mathbf{M}^{-1}=e^{[\mathrm{Ad}_{\mathbf{M}}\mathcal{A}]\theta}$（指数映射对共轭的协变性），可把 \cref{eq:ak-poe-space} 的每个 $e^{[\mathcal{S}_i]\theta_i}$ 改写。从 \cref{eq:ak-poe-space} 出发，反复插入 $\mathbf{M}\mathbf{M}^{-1}=\mathbf{I}$ 并把 $\mathbf{M}$ 一路右移过所有指数项：$\prod_i e^{[\mathcal{S}_i]\theta_i}\,\mathbf{M}=\mathbf{M}\prod_i\big(\mathbf{M}^{-1}e^{[\mathcal{S}_i]\theta_i}\mathbf{M}\big)=\mathbf{M}\prod_i e^{[\mathrm{Ad}_{\mathbf{M}^{-1}}\mathcal{S}_i]\theta_i}$，取 $\mathcal{B}_i=\mathrm{Ad}_{\mathbf{M}^{-1}}\mathcal{S}_i$ 即 \cref{eq:ak-poe-body}。物体形式的好处：每个 $\mathcal{B}_i$ 表达在 $\{n\}$ 中，构造下节的\emph{物体雅可比}时更直接。
\end{proof}

\begin{insight}[POE 三大优势：统一、少坐标系、贴合李群]
对照 DH，POE 的优势\cite{lynch2017modern}：(1) \textbf{统一处理转动/移动关节}——只是 $\boldsymbol{\omega}$ 是否为零之别，无需像 DH 那样分类讨论 $\theta_i$ vs $d_i$；(2) \textbf{只需两个坐标系}（$\{0\}$ 与零位 $\mathbf{M}$），免去逐连杆放系、免去标准/改进约定的歧义；(3) \textbf{与 \cref{ch:lie} 的李群机器无缝对接}——指数映射、伴随、旋量都是现成工具，雅可比（\cref{sec:ak-jacobian}）可直接由旋量读出。代价是要先求出零位 $\mathbf{M}$ 与各 $\mathcal{S}_i$，但这从 URDF（每关节的轴 + 安装位姿）可\emph{直接读出}，无额外建模负担。这就是为什么 Modern Robotics 把 POE 作为主线、本书亦然。
\end{insight}

\begin{pitfall}[陷阱：把 POE 的零位 $\mathbf{M}$ 当成“某个关节角下的位姿”]
$\mathbf{M}$ 是\emph{所有关节角为零}时的末端位姿，是一个\textbf{固定常量}，与运动无关。新手常误把它当成“初始构型”或随便某个位姿，导致旋量 $\mathcal{S}_i$（须在\emph{零位}测轴上点 $\mathbf{q}_i$）算错。正确流程：(1) 选定零位（通常取臂伸直或某规整构型）；(2) 在该零位下，对每个关节量出轴向 $\boldsymbol{\omega}_i$ 和轴上一点 $\mathbf{q}_i$（在 $\{0\}$ 下坐标），算 $\mathbf{v}_i=-\boldsymbol{\omega}_i\times\mathbf{q}_i$；(3) 量出该零位的末端位姿即 $\mathbf{M}$。三者都在\emph{同一个零位}下取，公式才自洽。
\end{pitfall}

\begin{practice}
\begin{enumerate}
  \item 对 2R 平面臂取零位为水平伸直，写出 $\mathcal{S}_1,\mathcal{S}_2$（两轴均沿 $z$，给出 $\mathbf{q}_1,\mathbf{q}_2$）与 $\mathbf{M}$，代入 \cref{eq:ak-poe-space} 验证与 §1.2 练习的 DH 结果一致。
  \item 证明 \cref{thm:ak-poe-body} 用到的恒等式 $\mathbf{M}e^{[\mathcal{A}]\theta}\mathbf{M}^{-1}=e^{[\mathrm{Ad}_{\mathbf{M}}\mathcal{A}]\theta}$（提示：对指数级数逐项用 $\mathbf{M}[\mathcal{A}]\mathbf{M}^{-1}=[\mathrm{Ad}_{\mathbf{M}}\mathcal{A}]$，\cref{eq:Ad}）。
  \item 一个移动关节沿 $\hat{\mathbf{z}}$ 方向。写出它的旋量 $\mathcal{S}$，并验证 $e^{[\mathcal{S}]d}$ 是沿 $z$ 平移 $d$ 的纯平移矩阵。
\end{enumerate}
\end{practice}

% ════════════════════════════════════════════════════════════════
\section{逆运动学：解析（闭式）解}
\label{sec:ak-ik-analytic}

\paragraph{本章的主旋律：“正易逆难”的不对称。}
读到这里，本章潜藏的核心张力终于浮出水面，值得在逆运动学开篇把它\emph{立为主线}。正运动学（FK）与逆运动学（IK）看似互逆，难度却\textbf{极不对称}：FK 是\emph{把链条正着连乘一遍}（\cref{eq:ak-poschain}），给定 $\mathbf{q}$ 必有唯一 $\mathbf{T}$、$O(n)$ 即得，再复杂的臂也不过多乘几个矩阵（\cref{thm:ak-poschain} 后的洞见已点出此事）；IK 却要\emph{反解一组非线性三角方程}，立刻撞上三道关卡——\emph{存在性}（目标可能在工作空间外、无解）、\emph{多解}（肘上/肘下、腕翻转，一般至多 8 组，\cref{note:ak-ik-multisol}）、\emph{闭式还是数值}（只有结构规整的臂才有闭式解，\cref{thm:ak-pieper}）。这道“正向走顺、反向走难”的鸿沟，是机械臂\emph{区别于}一般刚体（\cref{ch:rigid_body} 只需正向描述）的招牌难点，也是本章后半全部技术的\emph{同一个}根源：球腕之所以被设计出来，是为了\emph{驯服}逆解（本节）；奇异（\cref{sec:ak-singularity}）是这道鸿沟在\emph{速度层}的翻版——$\mathbf{J}$ 掉秩时反解 $\dot{\mathbf{q}}$ 同样爆炸，须用 DLS（\cref{eq:ak-dls}）正则化；冗余（\cref{sec:ak-redundancy}）则是鸿沟反向加宽——逆向多解多到成了\emph{无穷多}，反而可借零空间为我所用。Craig、Modern Robotics、Siciliano 诸经典无一例外都从这一不对称切入 IK\cite{craig2005introduction,lynch2017modern,siciliano2009robotics}。\textbf{抓住“正易逆难”这条主旋律，本章后半（IK／奇异／冗余／数值解）便不再是零散技巧，而是同一鸿沟的层层回响。}

\paragraph{动机：控制与规划都在任务空间下指令。}
我们想让末端\emph{到达某个位姿}（抓取点、插孔位姿），即给定 $\mathbf{T}_d$ 求关节角 $\mathbf{q}$ 使 $\mathbf{T}_{0n}(\mathbf{q})=\mathbf{T}_d$。这是\textbf{逆运动学}（inverse kinematics, IK）。与 FK 的“连乘即得”不同，IK 要解一组非线性方程，是机械臂运动学最有技术含量的部分。

\paragraph{反面：盲目数值迭代会卡在奇异、丢解。}
一个朴素想法是数值迭代（牛顿法 / 雅可比阻尼，见 \cref{sec:ak-singularity}）。但纯数值法有三个软肋：(1) 只收敛到\emph{一个}解，丢掉其它构型（肘上 / 肘下）；(2) 近奇异时雅可比病态、迭代发散；(3) 需初值、有局部极小。对\emph{结构规整}的臂（尤其带球腕的 6R），存在\textbf{闭式解析解}——一次算出全部解、无需迭代、无收敛问题，这是工业臂实时控制的首选。

\begin{theorem}[Pieper 准则：闭式解的充分条件]\label{thm:ak-pieper}
若 6 自由度串联臂满足下列任一结构条件，则其逆运动学存在\emph{闭式解析解}\cite{spong2006robot}：
\begin{enumerate}
  \item \textbf{三个相邻关节轴交于一点}（典型：球腕，腕三轴 $z_4,z_5,z_6$ 交于腕心）；或
  \item \textbf{三个相邻关节轴互相平行}。
\end{enumerate}
\end{theorem}
\begin{proof}[证明思路（位置--姿态解耦）]
以最常见的“前三轴定位 + 球腕”为例\cite{spong2006robot}。球腕三轴交于\emph{腕心} $\mathbf{p}_w$，故腕的三次转动\emph{不改变腕心位置}——腕心位置只由前三关节 $(q_1,q_2,q_3)$ 决定。给定目标 $\mathbf{T}_d=\big[\begin{smallmatrix}\mathbf{R}_d&\mathbf{p}_d\\\mathbf{0}^\top&1\end{smallmatrix}\big]$：
\begin{description}
  \item[第一步（解耦腕心）] 末端到腕心是沿末端 $z$ 轴的固定偏移 $d_6$，故腕心 $\mathbf{p}_w=\mathbf{p}_d-d_6\mathbf{R}_d\hat{\mathbf{z}}$ 可直接算出，\emph{只含已知量}。
  \item[第二步（解前三关节）] 由 $\mathbf{p}_w$ 反解 $(q_1,q_2,q_3)$——这是一个 3R 定位问题，用几何法（投影到平面、余弦定理）求闭式，一般得多组解（如肘上/肘下、肩左/肩右）。
  \item[第三步（解腕三关节）] 前三关节定后，$\mathbf{R}_{03}=\mathbf{R}_{03}(q_1,q_2,q_3)$ 已知，由 $\mathbf{R}_d=\mathbf{R}_{03}\mathbf{R}_{36}$ 得腕需提供的旋转 $\mathbf{R}_{36}=\mathbf{R}_{03}^\top\mathbf{R}_d$；把它当作 ZYZ（或 XYX，依腕轴次序）欧拉角分解即得 $(q_4,q_5,q_6)$，一般两组解（腕翻转）。
\end{description}
位置与姿态\emph{解耦}（第一步把二者分开）是闭式可解的关键；三轴交点保证“腕动不挪腕心”，三轴平行则保证另一种降维。
\end{proof}

\begin{insight}[球腕不是巧合，是“为可解性而设计”]
工业 6R 臂几乎都带球腕（腕三轴交于一点），这\emph{不是}制造巧合，而是\textbf{为了让 IK 有闭式解而刻意设计}的——Pieper 准则（\cref{thm:ak-pieper}）正是这一设计的理论依据。\cref{ch:arm-prac} 的 \texttt{arm\_dm} 臂“腰肩肘 + 球腕（$J_4,J_5,J_6$ 交点）”结构即如此。代价是球腕带来一类\emph{腕奇异}（$q_5=0$ 时 $z_4,z_6$ 共线，\cref{sec:ak-singularity}）；近年也有\emph{非球腕}臂（偏置腕）为避奇异/增工作空间而牺牲闭式可解性，转用数值或基于消元的半解析法（dialytic 消元，至多 16 解，\cref{sec:ak-ik-elimination}）。
\end{insight}

\begin{note}[多解、无解与工作空间]\label{note:ak-ik-multisol}
闭式 IK 自然暴露 IK 的本质特征\cite{spong2006robot}：
\begin{itemize}
  \item \textbf{多解}：6R 带球腕臂一般有最多 \textbf{8 组}解（前三关节 $2\times2=4$ 组肘/肩组合 $\times$ 腕 2 组翻转）。控制时须按\emph{最近解}（关节空间距当前构型最近）、\emph{避限位}、\emph{避奇异}择一。
  \item \textbf{无解}：目标在\emph{工作空间外}（超出臂展）或落在不可达的姿态域，则无实数解（几何法里表现为余弦定理的 $|\cos|>1$）。
  \item \textbf{工作空间}：可达点集（reachable workspace）与可任意定姿的点集（dexterous workspace）不同；后者是前者的真子集。规划时须确保目标在工作空间内（\cref{ch:arm-traj}）。
\end{itemize}
\end{note}

\paragraph{对照 P7：单腿解耦三步同此法。}
\cref{sec:qk-ik}（四足单腿逆运动学）用的正是同一“位置--姿态解耦 + 几何余弦定理”三步法：腿的 3R（髋 + 大腿 + 小腿）定足端位置，闭式解出三关节角，亦遇肘上/肘下式多解。本章把它\emph{推广}到 6R 带球腕的一般情形（Pieper 准则）；P7 是其在规整短链上的特例。读者可对照 \cref{sec:qk-ik} 体会“同一几何思想、不同臂型”的统一。

\begin{practice}
\begin{enumerate}
  \item 对 2R 平面臂，用余弦定理由目标位置 $(x,y)$ 解出 $\theta_2=\pm\arccos\frac{x^2+y^2-a_1^2-a_2^2}{2a_1a_2}$（两解＝肘上/肘下），再解 $\theta_1$。讨论何时无解。
  \item 验证 Pieper 第一步：若末端到腕心沿末端 $z$ 偏移 $d_6$，则 $\mathbf{p}_w=\mathbf{p}_d-d_6\mathbf{R}_d\hat{\mathbf{z}}$。为什么腕的三次转动不改变 $\mathbf{p}_w$？
  \item 一个 6R 臂在某目标位姿处恰好有 8 组 IK 解。若该臂某关节限位收紧使肘只能“上”，问可行解减到几组？
\end{enumerate}
\end{practice}

% ----------------------------------------------
\subsection{螺旋式逆运动学：Paden--Kahan 子问题}
\label{sec:ak-ik-paden-kahan}

\paragraph{动机：让逆解\emph{长在} POE 的李群机器上。}
Pieper 准则（\cref{thm:ak-pieper}）的解耦三步固然给出闭式解，却\emph{换了一套语言}：把刚建好的 POE/旋量暂时搁下，改用腕心坐标、再把 $\mathbf{R}_{36}$ 当 ZYZ 欧拉角分解去凑解。这既要在旋量与欧拉角之间来回切换，又对欧拉角的分支选择、万向锁奇异格外敏感。Paden 与 Kahan 给出另一条路\cite{paden1986thesis,murray1994mathematical}：\textbf{直接在指数积方程 $e^{[\mathcal{S}_1]\theta_1}\cdots e^{[\mathcal{S}_n]\theta_n}\mathbf{M}=g_d$ 上做几何反解}——把整条方程\emph{拆成几个有现成闭式解、几何意义清澈、数值稳定的标准子问题}，逐个解出关节角。如此逆运动学\emph{完全复用} \cref{ch:lie} 与 \cref{sec:ak-poe} 已备好的 $\exp([\mathcal{S}]\theta)$、伴随、旋量机器，不再另起炉灶。这是把“正易逆难”鸿沟（\cref{sec:ak-ik-analytic} 开篇）用\emph{李群语言}驯服的又一招。

\begin{note}[本小节局部记号]\label{note:ak-pk-notation}
本小节沿用 Paden--Kahan 文献记号：$\mathbf{p},\mathbf{q}\in\mathbb{R}^3$ 记空间中两\emph{点}、$\mathbf{r}$ 记轴上一点、$\boldsymbol{\omega}$ 记轴向单位向量，\emph{均非}关节向量 $\mathbf{q}$（\cref{note:ak-notation}）。零节距（zero-pitch）单位旋量 $\mathcal{S}=[\boldsymbol{\omega};\,\mathbf{v}]$ 即纯转动关节旋量（\cref{def:ak-screw}，$\mathbf{v}=-\boldsymbol{\omega}\times\mathbf{r}$），$e^{[\mathcal{S}]\theta}$ 作用于（齐次）点。$\mathbf{u}'=\mathbf{u}-\boldsymbol{\omega}\boldsymbol{\omega}^\top\mathbf{u}$ 记 $\mathbf{u}$ 在垂直 $\boldsymbol{\omega}$ 平面上的投影。
\end{note}

\paragraph{核心思想：两把“约化剪刀”。}
把含多个指数的方程化成单个子问题，靠两个反复使用的技巧\cite{murray1994mathematical}：
\begin{enumerate}
  \item \textbf{作用到轴上点（消去末关节）}：若点 $\mathbf{p}$ 落在转动旋量 $\mathcal{S}_k$ 的轴上，则 $e^{[\mathcal{S}_k]\theta_k}\mathbf{p}=\mathbf{p}$，把方程两端作用到 $\mathbf{p}$ 即\emph{消去} $\theta_k$。例：$e^{[\mathcal{S}_1]\theta_1}e^{[\mathcal{S}_2]\theta_2}e^{[\mathcal{S}_3]\theta_3}=g$ 作用到 $\mathcal{S}_3$ 轴上点 $\mathbf{p}$ 得 $e^{[\mathcal{S}_1]\theta_1}e^{[\mathcal{S}_2]\theta_2}\mathbf{p}=g\mathbf{p}$（子问题 2）。
  \item \textbf{减点取模（用刚体保距消去转动）}：刚体运动保持点距，故对两端减去某点 $\mathbf{q}$ 再取模可消去“绕 $\mathbf{q}$ 的转动”。例：上式中若 $\mathcal{S}_1,\mathcal{S}_2$ 交于 $\mathbf{q}$，对 $\mathcal{S}_3$ 轴外一点 $\mathbf{p}$ 有 $\|g\mathbf{p}-\mathbf{q}\|=\|e^{[\mathcal{S}_1]\theta_1}e^{[\mathcal{S}_2]\theta_2}(e^{[\mathcal{S}_3]\theta_3}\mathbf{p}-\mathbf{q})\|=\|e^{[\mathcal{S}_3]\theta_3}\mathbf{p}-\mathbf{q}\|$（子问题 3）。
\end{enumerate}
下面三个子问题是这些约化的\emph{终点}——它们都有几何闭式解。

\begin{figure}[ht]
\centering
\begin{tikzpicture}[scale=0.95,>=Stealth,font=\scriptsize]
% Panel 1: subproblem 1 — point on a circle about one axis
\begin{scope}[shift={(0,0)}]
  \draw[->,thick,gray!50!black] (-0.1,-1.05)--(-0.1,1.35) node[above]{$\boldsymbol{\omega}$};
  \draw[main!70!black,thick] (0,0) ellipse (1 and 0.3);
  \fill (1,0) circle (1.3pt) node[right]{$\mathbf{p}$};
  \fill (-0.45,0.27) circle (1.3pt) node[above left=-2pt]{$\mathbf{q}$};
  \draw[->,second!70!black,thick] (0,0)--(1,0);
  \draw[->,second!70!black,thick] (0,0)--(-0.45,0.27);
  \node at (0.28,0.16){$\theta$};
  \node at (0,-1.5){子问题 1：绕单轴 $\mathbf{p}\!\to\!\mathbf{q}$};
\end{scope}
% Panel 2: subproblem 2 — two circles intersect at c
\begin{scope}[shift={(4.0,0)}]
  \draw[main!70!black,thick] (-0.5,0.05) ellipse (0.62 and 0.92);
  \draw[third!70!black,thick] (0.5,0.05) ellipse (0.62 and 0.92);
  \fill (-0.5,0.97) circle (1.3pt) node[above]{$\mathbf{p}$};
  \fill (0.5,0.97) circle (1.3pt) node[above]{$\mathbf{q}$};
  \fill (0,-0.62) circle (1.3pt) node[below]{$\mathbf{c}$};
  \node at (0,-1.5){子问题 2：两圆交于 $\mathbf{c}$};
\end{scope}
% Panel 3: subproblem 3 — circle meets sphere of radius delta about q
\begin{scope}[shift={(7.8,0)}]
  \draw[main!70!black,thick] (0,0) ellipse (0.95 and 0.3);
  \fill (0.95,0) circle (1.3pt) node[right]{$\mathbf{p}$};
  \fill (-0.25,1.0) circle (1.3pt) node[above]{$\mathbf{q}$};
  \draw[dashed,second!70!black] (-0.25,1.0) circle (0.82);
  \draw[->,second!70!black] (-0.25,1.0)--(0.52,0.72) node[midway,above=-1pt]{$\delta$};
  \node at (0,-1.5){子问题 3：转到距 $\mathbf{q}$ 为 $\delta$};
\end{scope}
\end{tikzpicture}
\caption{Paden--Kahan 三子问题的几何。子问题 1：把 $\mathbf{p}$ 绕单轴转到 $\mathbf{q}$（$\mathbf{q}$ 须落在 $\mathbf{p}$ 所画的圆上）；子问题 2：$\mathbf{p}$ 绕轴 2、$\mathbf{q}$ 绕轴 1 各画一圆，过渡点 $\mathbf{c}$ 是两圆交点（至多两交点＝两组解）；子问题 3：把 $\mathbf{p}$ 绕轴转到与 $\mathbf{q}$ 相距 $\delta$（圆与以 $\mathbf{q}$ 为心、$\delta$ 为半径的球面相交）。三者各对应 $0/1/2$ 解的几何相交计数。}
\label{fig:ak-paden-kahan}
\end{figure}

\begin{theorem}[Paden--Kahan 子问题 1：绕单轴旋转]\label{thm:ak-pk-sub1}
设 $\mathcal{S}$ 为零节距单位旋量、轴向 $\boldsymbol{\omega}$、轴上一点 $\mathbf{r}$；$\mathbf{p},\mathbf{q}\in\mathbb{R}^3$。求 $\theta$ 使\cite{murray1994mathematical}
\begin{equation}
e^{[\mathcal{S}]\theta}\mathbf{p}=\mathbf{q}.
\label{eq:ak-pk-sub1}
\end{equation}
令 $\mathbf{u}=\mathbf{p}-\mathbf{r}$、$\mathbf{v}=\mathbf{q}-\mathbf{r}$ 及其垂直投影 $\mathbf{u}'=\mathbf{u}-\boldsymbol{\omega}\boldsymbol{\omega}^\top\mathbf{u}$、$\mathbf{v}'=\mathbf{v}-\boldsymbol{\omega}\boldsymbol{\omega}^\top\mathbf{v}$。则 \cref{eq:ak-pk-sub1} 有解\emph{当且仅当}
\begin{equation}
\boldsymbol{\omega}^\top\mathbf{u}=\boldsymbol{\omega}^\top\mathbf{v}\quad\text{且}\quad\|\mathbf{u}'\|=\|\mathbf{v}'\|;
\label{eq:ak-pk-sub1-cond}
\end{equation}
此时（$\mathbf{u}'\neq\mathbf{0}$）唯一解为
\begin{equation}
\theta=\operatorname{atan2}\!\big(\boldsymbol{\omega}^\top(\mathbf{u}'\times\mathbf{v}'),\ \mathbf{u}'^\top\mathbf{v}'\big).
\label{eq:ak-pk-sub1-sol}
\end{equation}
若 $\mathbf{u}'=\mathbf{0}$（$\mathbf{p},\mathbf{q}$ 都在轴上），则 $\theta$ 任意，无穷多解。
\end{theorem}
\begin{proof}
绕轴转动保持点的\emph{沿轴高度}与\emph{到轴距离}，故 $\mathbf{p}$ 能绕 $\mathcal{S}$ 转到 $\mathbf{q}$ 当且仅当二者沿轴分量相等（$\boldsymbol{\omega}^\top\mathbf{u}=\boldsymbol{\omega}^\top\mathbf{v}$）且到轴距离相等（$\|\mathbf{u}'\|=\|\mathbf{v}'\|$），即可解条件 \cref{eq:ak-pk-sub1-cond}。因 $\mathbf{r}$ 在轴上 $e^{[\mathcal{S}]\theta}\mathbf{r}=\mathbf{r}$，\cref{eq:ak-pk-sub1} 等价于绕过原点轴 $\boldsymbol{\omega}$ 的纯旋转 $e^{\boldsymbol{\omega}^\wedge\theta}\mathbf{u}=\mathbf{v}$（\cref{ch:lie} 的 $\SO(3)$ 指数映射）。转动只发生在垂直 $\boldsymbol{\omega}$ 的平面内，对投影 $\mathbf{u}',\mathbf{v}'$ 有 $\boldsymbol{\omega}^\top(\mathbf{u}'\times\mathbf{v}')=\|\mathbf{u}'\|\|\mathbf{v}'\|\sin\theta$、$\mathbf{u}'^\top\mathbf{v}'=\|\mathbf{u}'\|\|\mathbf{v}'\|\cos\theta$，二者经 $\operatorname{atan2}$ 给出 \cref{eq:ak-pk-sub1-sol}，自动落在正确象限（$\operatorname{atan2}$ 而非 $\arccos$ 主值——这是数值稳定的关键）。
\end{proof}

\begin{theorem}[Paden--Kahan 子问题 2：绕两相交轴顺次旋转]\label{thm:ak-pk-sub2}
设 $\mathcal{S}_1,\mathcal{S}_2$ 为零节距单位旋量、轴向 $\boldsymbol{\omega}_1,\boldsymbol{\omega}_2$ \emph{相交}于点 $\mathbf{r}$（$\boldsymbol{\omega}_1\times\boldsymbol{\omega}_2\neq\mathbf{0}$）；$\mathbf{p},\mathbf{q}\in\mathbb{R}^3$。求 $\theta_1,\theta_2$ 使\cite{murray1994mathematical}
\begin{equation}
e^{[\mathcal{S}_1]\theta_1}e^{[\mathcal{S}_2]\theta_2}\mathbf{p}=\mathbf{q}.
\label{eq:ak-pk-sub2}
\end{equation}
取过渡点 $\mathbf{c}$ 满足 $e^{[\mathcal{S}_2]\theta_2}\mathbf{p}=\mathbf{c}=e^{-[\mathcal{S}_1]\theta_1}\mathbf{q}$，并把 $\mathbf{z}=\mathbf{c}-\mathbf{r}$ 按基 $\{\boldsymbol{\omega}_1,\boldsymbol{\omega}_2,\boldsymbol{\omega}_1\times\boldsymbol{\omega}_2\}$ 展开 $\mathbf{z}=\alpha\boldsymbol{\omega}_1+\beta\boldsymbol{\omega}_2+\gamma(\boldsymbol{\omega}_1\times\boldsymbol{\omega}_2)$，其中（$\mathbf{u}=\mathbf{p}-\mathbf{r}$、$\mathbf{v}=\mathbf{q}-\mathbf{r}$）
\begin{equation}
\alpha=\frac{(\boldsymbol{\omega}_1^\top\boldsymbol{\omega}_2)\,\boldsymbol{\omega}_2^\top\mathbf{u}-\boldsymbol{\omega}_1^\top\mathbf{v}}{(\boldsymbol{\omega}_1^\top\boldsymbol{\omega}_2)^2-1},\quad
\beta=\frac{(\boldsymbol{\omega}_1^\top\boldsymbol{\omega}_2)\,\boldsymbol{\omega}_1^\top\mathbf{v}-\boldsymbol{\omega}_2^\top\mathbf{u}}{(\boldsymbol{\omega}_1^\top\boldsymbol{\omega}_2)^2-1},\quad
\gamma^2=\frac{\|\mathbf{u}\|^2-\alpha^2-\beta^2-2\alpha\beta\,\boldsymbol{\omega}_1^\top\boldsymbol{\omega}_2}{\|\boldsymbol{\omega}_1\times\boldsymbol{\omega}_2\|^2}.
\label{eq:ak-pk-sub2-gamma}
\end{equation}
$\gamma^2<0/=0/>0$ 分别给 $0/1/2$ 个 $\mathbf{c}$；对每个 $\mathbf{c}$，用子问题 1 解 $e^{[\mathcal{S}_2]\theta_2}\mathbf{p}=\mathbf{c}$ 得 $\theta_2$、$e^{-[\mathcal{S}_1]\theta_1}\mathbf{q}=\mathbf{c}$ 得 $\theta_1$。
\end{theorem}
\begin{proof}
$\mathbf{p}$ 先绕 $\mathcal{S}_2$ 转到 $\mathbf{c}$、$\mathbf{c}$ 再绕 $\mathcal{S}_1$ 转到 $\mathbf{q}$，故 $\mathbf{c}$ 同时落在“$\mathbf{p}$ 绕轴 2 画的圆”与“$\mathbf{q}$ 绕轴 1 画的圆”上，二圆交点即 $\mathbf{c}$（\cref{fig:ak-paden-kahan} 中栏）。两轴交于 $\mathbf{r}$，平移使 $\mathbf{r}$ 为原点得 $e^{\boldsymbol{\omega}_2^\wedge\theta_2}\mathbf{u}=\mathbf{z}=e^{-\boldsymbol{\omega}_1^\wedge\theta_1}\mathbf{v}$。分别点乘 $\boldsymbol{\omega}_2,\boldsymbol{\omega}_1$（旋转不改变沿自身轴投影）得 $\boldsymbol{\omega}_2^\top\mathbf{u}=\boldsymbol{\omega}_2^\top\mathbf{z}$、$\boldsymbol{\omega}_1^\top\mathbf{v}=\boldsymbol{\omega}_1^\top\mathbf{z}$，再加等距 $\|\mathbf{u}\|=\|\mathbf{z}\|=\|\mathbf{v}\|$。因 $\boldsymbol{\omega}_1,\boldsymbol{\omega}_2$ 线性无关，$\{\boldsymbol{\omega}_1,\boldsymbol{\omega}_2,\boldsymbol{\omega}_1\times\boldsymbol{\omega}_2\}$ 是 $\mathbb{R}^3$ 一组基；把 $\mathbf{z}$ 的展开式代入前两个投影方程得关于 $\alpha,\beta$ 的二元线性方程组，解得 \cref{eq:ak-pk-sub2-gamma} 前两式；用 $\|\mathbf{z}\|^2=\|\mathbf{u}\|^2$ 解出 $\gamma^2$ 得第三式。$\gamma^2$ 的符号即两圆相离/相切/相交，对应 $0/1/2$ 解。
\end{proof}

\begin{theorem}[Paden--Kahan 子问题 3：转到给定距离]\label{thm:ak-pk-sub3}
设 $\mathcal{S}$ 为零节距单位旋量、轴向 $\boldsymbol{\omega}$、轴上点 $\mathbf{r}$；$\mathbf{p},\mathbf{q}\in\mathbb{R}^3$，$\delta>0$。求 $\theta$ 使转后 $\mathbf{p}$ 与 $\mathbf{q}$ 相距 $\delta$\cite{murray1994mathematical}：
\begin{equation}
\big\|\mathbf{q}-e^{[\mathcal{S}]\theta}\mathbf{p}\big\|=\delta.
\label{eq:ak-pk-sub3}
\end{equation}
令 $\mathbf{u}=\mathbf{p}-\mathbf{r}$、$\mathbf{v}=\mathbf{q}-\mathbf{r}$ 及投影 $\mathbf{u}',\mathbf{v}'$，把距离投到垂面 $\delta'^2=\delta^2-\big|\boldsymbol{\omega}^\top(\mathbf{p}-\mathbf{q})\big|^2$，并记
\begin{equation}
\theta_0=\operatorname{atan2}\!\big(\boldsymbol{\omega}^\top(\mathbf{u}'\times\mathbf{v}'),\ \mathbf{u}'^\top\mathbf{v}'\big).
\label{eq:ak-pk-sub3-theta0}
\end{equation}
则由\emph{余弦定理}
\begin{equation}
\boxed{\ \theta=\theta_0\pm\cos^{-1}\!\left(\frac{\|\mathbf{u}'\|^2+\|\mathbf{v}'\|^2-\delta'^2}{2\,\|\mathbf{u}'\|\,\|\mathbf{v}'\|}\right).\ }
\label{eq:ak-pk-sub3-cos}
\end{equation}
据 $\|\mathbf{u}'\|$ 圆与半径 $\delta'$ 圆的交点数，\cref{eq:ak-pk-sub3} 有 $0/1/2$ 解。
\end{theorem}
\begin{proof}
转动保持沿轴分量，故 $\mathbf{p},\mathbf{q}$ 沿 $\boldsymbol{\omega}$ 的距离分量 $\big|\boldsymbol{\omega}^\top(\mathbf{p}-\mathbf{q})\big|$ 与 $\theta$ 无关，可先从 $\delta^2$ 扣除，余下垂面内距离平方 $\delta'^2$。在垂面内，$\mathbf{u}'$ 转过 $\theta$ 后的端点与 $\mathbf{v}'$ 端点相距 $\delta'$；轴心、$e^{\boldsymbol{\omega}^\wedge\theta}\mathbf{u}'$ 端、$\mathbf{v}'$ 端三点构成三角形，两边长 $\|\mathbf{u}'\|,\|\mathbf{v}'\|$、对边 $\delta'$、夹角 $\phi=\theta_0-\theta$（$\theta_0$ 是 $\mathbf{u}',\mathbf{v}'$ 间的夹角，\cref{eq:ak-pk-sub3-theta0}）。\textbf{余弦定理} $\|\mathbf{u}'\|^2+\|\mathbf{v}'\|^2-2\|\mathbf{u}'\|\|\mathbf{v}'\|\cos\phi=\delta'^2$ 解出 $\phi$，$\theta=\theta_0-\phi$ 即 \cref{eq:ak-pk-sub3-cos}，$\pm$ 对应两交点。
\end{proof}

\begin{note}[更多子问题：本组并不“穷尽”]\label{note:ak-pk-more}
\cref{thm:ak-pk-sub1,thm:ak-pk-sub2,thm:ak-pk-sub3} 是\emph{最常用}的三个，但非全部\cite{murray1994mathematical}：还有“绕两平行轴转到点”“沿移动副+转动副组合”等变体，处理含移动关节或特殊轴序的臂（习题里再展开）。它们共享同一哲学：\textbf{每个子问题对应一类“圆/球/直线相交”的初等几何，闭式可解、解数＝相交计数、用 $\operatorname{atan2}$ 数值稳健}。给一台新臂做 Paden--Kahan 逆解，关键工夫在\emph{选对作用点}把方程约化到已知子问题，而非求解子问题本身。
\end{note}

\begin{example}[肘式 6R 机械臂的螺旋式逆运动学（四步分解）]\label{exam:ak-elbow-ik}
\textbf{肘式}（elbow）机械臂是带球腕的 6R 臂：前三轴 $\mathcal{S}_1,\mathcal{S}_2,\mathcal{S}_3$ 定位（腰、肩、肘），后三轴 $\mathcal{S}_4,\mathcal{S}_5,\mathcal{S}_6$ 交于腕心 $\mathbf{p}_w$ 定姿（球腕）——工业臂的典型结构（\cref{thm:ak-pieper}），最贴合 Paden--Kahan 分解\cite{murray1994mathematical}。给定目标 $g_d\in\SE(3)$，要解
\begin{equation}
g_{st}(\boldsymbol{\theta})=e^{[\mathcal{S}_1]\theta_1}\cdots e^{[\mathcal{S}_6]\theta_6}\,\mathbf{M}=g_d
\;\Longrightarrow\;
e^{[\mathcal{S}_1]\theta_1}\cdots e^{[\mathcal{S}_6]\theta_6}=g_d\mathbf{M}^{-1}=:g_1.
\label{eq:ak-elbow-g1}
\end{equation}
分四步逐组解出 $\theta_1,\dots,\theta_6$，每步恰好约化成一个子问题。

\textbf{第一步（肘角 $\theta_3$，子问题 3）.} 把 \cref{eq:ak-elbow-g1} 两端作用到腕心 $\mathbf{p}_w$。因 $\mathbf{p}_w$ 在三条腕轴上，$e^{[\mathcal{S}_4]\theta_4}e^{[\mathcal{S}_5]\theta_5}e^{[\mathcal{S}_6]\theta_6}\mathbf{p}_w=\mathbf{p}_w$，于是腕三角\emph{被消去}：
\begin{equation}
e^{[\mathcal{S}_1]\theta_1}e^{[\mathcal{S}_2]\theta_2}e^{[\mathcal{S}_3]\theta_3}\mathbf{p}_w=g_1\mathbf{p}_w.
\label{eq:ak-elbow-step1}
\end{equation}
再两端减前两轴交点 $\mathbf{p}_b$（在 $\mathcal{S}_1,\mathcal{S}_2$ 轴上）并取模——$e^{[\mathcal{S}_1]\theta_1}e^{[\mathcal{S}_2]\theta_2}$ 是绕 $\mathbf{p}_b$ 的转动、保持到 $\mathbf{p}_b$ 的距离：
\begin{equation}
\big\|e^{[\mathcal{S}_3]\theta_3}\mathbf{p}_w-\mathbf{p}_b\big\|=\big\|g_1\mathbf{p}_w-\mathbf{p}_b\big\|.
\label{eq:ak-elbow-norm}
\end{equation}
这正是\textbf{子问题 3}（\cref{thm:ak-pk-sub3}：$\mathbf{p}=\mathbf{p}_w$、$\mathbf{q}=\mathbf{p}_b$、$\delta=\|g_1\mathbf{p}_w-\mathbf{p}_b\|$），解出 $\theta_3$（至多 2 解＝肘上/肘下）。

\textbf{第二步（基座两角 $\theta_1,\theta_2$，子问题 2）.} $\theta_3$ 已知，\cref{eq:ak-elbow-step1} 成为 $e^{[\mathcal{S}_1]\theta_1}e^{[\mathcal{S}_2]\theta_2}\big(e^{[\mathcal{S}_3]\theta_3}\mathbf{p}_w\big)=g_1\mathbf{p}_w$，即\textbf{子问题 2}（\cref{thm:ak-pk-sub2}：$\mathbf{p}=e^{[\mathcal{S}_3]\theta_3}\mathbf{p}_w$、$\mathbf{q}=g_1\mathbf{p}_w$，$\mathcal{S}_1,\mathcal{S}_2$ 交于 $\mathbf{p}_b$），解出 $\theta_1,\theta_2$（至多 2 解＝肩左/肩右）。

\textbf{第三步（腕两角 $\theta_4,\theta_5$，子问题 2）.} 前三角已知，剩余腕部满足 $e^{[\mathcal{S}_4]\theta_4}e^{[\mathcal{S}_5]\theta_5}e^{[\mathcal{S}_6]\theta_6}=e^{-[\mathcal{S}_3]\theta_3}e^{-[\mathcal{S}_2]\theta_2}e^{-[\mathcal{S}_1]\theta_1}g_1=:g_2$。取点 $\mathbf{p}$ 在 $\mathcal{S}_6$ 轴上但\emph{不}在 $\mathcal{S}_4,\mathcal{S}_5$ 轴上，$e^{[\mathcal{S}_6]\theta_6}\mathbf{p}=\mathbf{p}$，得 $e^{[\mathcal{S}_4]\theta_4}e^{[\mathcal{S}_5]\theta_5}\mathbf{p}=g_2\mathbf{p}$——又是\textbf{子问题 2}（$\mathcal{S}_4,\mathcal{S}_5$ 交于腕心），解出 $\theta_4,\theta_5$（至多 2 解＝腕翻转）。

\textbf{第四步（末腕角 $\theta_6$，子问题 1）.} 唯一未知 $\theta_6$。取 $\mathbf{p}$ \emph{不}在 $\mathcal{S}_6$ 轴上，整理腕方程为 $e^{[\mathcal{S}_6]\theta_6}\mathbf{p}=e^{-[\mathcal{S}_5]\theta_5}e^{-[\mathcal{S}_4]\theta_4}g_2\mathbf{p}=:\mathbf{q}$，即\textbf{子问题 1}（\cref{thm:ak-pk-sub1}），解出 $\theta_6$。

\textbf{解数核对.} 全程解数＝第一步 $2$（肘）$\times$ 第二步 $2$（肩）$\times$ 第三步 $2$（腕翻转）$=\mathbf{8}$ 组——与 \cref{note:ak-ik-multisol} 的“6R 带球腕至多 8 解”从\emph{子问题角度}重数了一遍，且每组解都由 $\operatorname{atan2}$/余弦定理\emph{一次性算出}，无迭代、无欧拉角奇异。
\end{example}

\begin{insight}[同一闭式解，两副面孔：Pieper 的欧拉角 vs Paden--Kahan 的子问题]\label{ins:ak-pieper-vs-pk}
对照 \cref{thm:ak-pieper}：Pieper 把 IK 解耦成“腕心位置 $+$ $\mathbf{R}_{36}$ 的 ZYZ 欧拉角”，本质是\emph{在坐标分量上}凑解；Paden--Kahan 把\emph{同一台臂的同一组}解，改写成三类\emph{坐标无关}的几何子问题（绕轴转到点 / 两圆交 / 圆与球交），全程在 $\exp([\mathcal{S}]\theta)$ 上操作。两者解数（8 组）、解集\emph{完全相同}，差别在方法论：Paden--Kahan (1) 复用 POE 机器、不切换到欧拉角；(2) 每子问题几何清晰、数值稳健（$\operatorname{atan2}$ 而非 $\arccos$ 主值）；(3) 易推广到非标准轴序。代价是要\emph{识别}哪一步约化成哪个子问题（选对作用点：轴上点消元、相交点取模）。这印证 \cref{sec:ak-ik-analytic} 开篇的主旋律——\textbf{球腕的可解性既能被坐标法（Pieper）也能被李群法（Paden--Kahan）兑现，后者把逆解\emph{几何化}、与本书 POE 主线一脉相承。}
\end{insight}

\begin{pitfall}[陷阱：作用点选错，约化不成子问题]
Paden--Kahan 的全部技巧在\emph{选对作用点}。两个高频错误：(1) 把方程作用到\emph{不在}目标关节轴上的点，$e^{[\mathcal{S}_k]\theta_k}\mathbf{p}\neq\mathbf{p}$，消不掉 $\theta_k$，约化失败；(2) 用“减点取模”时，所减的点\emph{不是}两轴交点，刚体保距性不成立，得不到子问题 3。正确流程：先在臂上\emph{标出}各关节轴的交点（腕心 $\mathbf{p}_w$、基座交点 $\mathbf{p}_b$）与轴上代表点，再按“要消哪个 $\theta$ 就作用到哪条轴上点”逐步约化。另注意：子问题 2 要求两轴\emph{真正相交}（$\boldsymbol{\omega}_1\times\boldsymbol{\omega}_2\neq\mathbf{0}$ 且共点）；轴\emph{异面}（既不交也不平行）时本子问题不适用，须换更一般的子问题或转下节消元法。
\end{pitfall}

\begin{practice}
\begin{enumerate}
  \item 用子问题 1 重解 2R 平面臂第二关节：取 $\mathcal{S}_2$（沿 $z$、过肘）、$\mathbf{p}=$ 零位末端、$\mathbf{q}=$ 目标末端绕肩转回，验证与 §1.4 余弦定理两解一致。
  \item 证明子问题 3 的 $\delta'^2=\delta^2-|\boldsymbol{\omega}^\top(\mathbf{p}-\mathbf{q})|^2$：为什么沿轴分量可先扣除？（提示：转动不改变沿 $\boldsymbol{\omega}$ 的坐标。）
  \item 对 \cref{exam:ak-elbow-ik} 第一步，说明为何要\emph{先}作用到 $\mathbf{p}_w$（消腕三角）\emph{再}减 $\mathbf{p}_b$ 取模（消肩腰两角），两步顺序能否颠倒？
  \item 若某 6R 臂腕三轴\emph{不}交于一点（偏置腕），\cref{exam:ak-elbow-ik} 第一步的“$e^{[\mathcal{S}_4]\theta_4}e^{[\mathcal{S}_5]\theta_5}e^{[\mathcal{S}_6]\theta_6}\mathbf{p}_w=\mathbf{p}_w$”还成立吗？这预示了下节为何需要消元法。
\end{enumerate}
\end{practice}

% ----------------------------------------------
\subsection{通用 6R 逆运动学：消元法与解数上界 16}
\label{sec:ak-ik-elimination}

\paragraph{动机：当“没有相交轴”可借力时。}
Paden--Kahan 子问题（\cref{sec:ak-ik-paden-kahan}）与 Pieper 准则（\cref{thm:ak-pieper}）都\emph{依赖结构}——要有三轴交点、或相交/平行轴可供“作用到轴上点”消元。但一台\emph{通用} 6R 臂（六轴位置、方向\emph{一般}：偏置腕、为避奇异/增工作空间而故意打散轴系，\cref{exam:ak-elbow-ik} 练习 4 已预告）没有这种特殊结构，子问题\emph{无从下手}。此时只剩\emph{纯代数}一条路：把运动学方程化成多项式方程组，用代数几何的\textbf{消元理论}（elimination theory）系统地消去其中 $n-1$ 个变量，归结为\emph{一个}单变量高次多项式，求根即得全部解。本节按 Manocha--Canny、Raghavan--Roth、Lee--Liang 的经典工作\cite{manocha1994realtime,raghavan1993inverse,lee1988displacement,murray1994mathematical} 给出这条路线，并收获一个漂亮而尖锐的结论：通用 6R 逆解\emph{至多 16 组}。

\paragraph{核心工具：dialytic（戴利）消元。}
设有 $n$ 个多项式、$n$ 个变量，要消去其中 $n-1$ 个、得到关于剩余一个变量的单变量多项式。\textbf{dialytic 消元}的思路是：把方程组看成关于“单项式向量”的\emph{线性}方程组，\emph{不断用单项式去乘已有方程、生成新方程}，直到方程个数等于出现的单项式个数；此时系数矩阵为方阵，方程组有非零解（单项式向量非零）\emph{当且仅当}系数矩阵\emph{行列式为零}——这个行列式即所求的单变量多项式。它是“暴力但通用”的过程：只用到方程的极少结构，故对任意 6R 臂都适用（不像子问题法要求特殊轴系）。先用一个纯代数的暖身例熟悉机理。

\begin{example}[三元多项式的 dialytic 消元（暖身）]\label{exam:ak-dialytic}
设三个关于 $x_1,x_2,x_3$ 的多项式 $f_1,f_2,f_3$ 取如下（6 单项式）形式\cite{murray1994mathematical}：
\begin{equation}
\begin{bmatrix}f_1\\f_2\\f_3\end{bmatrix}
=\begin{bmatrix}a_1&b_1&c_1&d_1&e_1&g_1\\a_2&b_2&c_2&d_2&e_2&g_2\\a_3&b_3&c_3&d_3&e_3&g_3\end{bmatrix}
\begin{bmatrix}x_1^2\\x_1x_2\\x_1x_3\\x_2^2\\x_3\\1\end{bmatrix}=\mathbf{0},
\label{eq:ak-dialytic-f}
\end{equation}
$a_i,\dots,g_i\in\mathbb{R}$。欲消去 $x_2,x_3$、求关于 $x_1$ 的公共根多项式，分两步。

\textbf{第一步（视为 $x_2,x_3$ 的方程、系数含 $x_1$）.} 把 \cref{eq:ak-dialytic-f} 按 $x_2,x_3$ 的单项式 $\{x_2^2,x_3,x_2,1\}$ 重组（$x_1$ 当参数）：
\begin{equation}
\begin{bmatrix}f_1\\f_2\\f_3\end{bmatrix}
=\begin{bmatrix}d_1&e_1+c_1x_1&b_1x_1&a_1x_1^2+g_1\\ d_2&e_2+c_2x_1&b_2x_1&a_2x_1^2+g_2\\ d_3&e_3+c_3x_1&b_3x_1&a_3x_1^2+g_3\end{bmatrix}
\begin{bmatrix}x_2^2\\x_3\\x_2\\1\end{bmatrix}=\mathbf{0}.
\label{eq:ak-dialytic-regroup}
\end{equation}
这是 $3$ 个方程、$4$ 个单项式——方程不够，系数矩阵非方。

\textbf{第二步（乘 $x_2$ 生成新方程、配平到 $6\times6$）.} 把 $f_1,f_2,f_3$ 各乘 $x_2$，引入新单项式 $\{x_2^3,x_2x_3,x_2^2,x_2\}$；与原方程合并，单项式集闭合为 $6$ 个 $\{x_2^3,x_2x_3,x_2^2,x_3,x_2,1\}$，方程也凑足 $6$ 个：
\begin{equation}
A(x_1)\begin{bmatrix}x_2^3\\x_2x_3\\x_2^2\\x_3\\x_2\\1\end{bmatrix}=\mathbf{0},\quad
A(x_1)=\begin{bmatrix}
0&0&d_i&e_i+c_ix_1&b_ix_1&a_ix_1^2+g_i\\[-1pt]
d_i&e_i+c_ix_1&b_ix_1&0&a_ix_1^2+g_i&0
\end{bmatrix}_{i=1,2,3},
\label{eq:ak-dialytic-A}
\end{equation}
其中上三行（$i=1,2,3$）来自原 $f_i$、下三行来自 $x_2f_i$，$A(x_1)\in\mathbb{R}^{6\times6}$。单项式向量非零，故须
\begin{equation}
\det A(x_1)=0.
\label{eq:ak-dialytic-det}
\end{equation}
这是关于 $x_1$ 的\emph{单变量}多项式（一般\emph{六次}）；求其根得 $x_1$，回代 \cref{eq:ak-dialytic-A} 的零空间（最后一项归一）即读出对应的 $x_2,x_3$。\textbf{消元完成：三元三式 $\to$ 一元六次。}
\end{example}

\begin{insight}[dialytic 消元的“配平”哲学]\label{ins:ak-dialytic}
\cref{exam:ak-dialytic} 的关键动作是\emph{故意制造冗余}：原本 3 式不足以定 4 单项式，便用 $x_2$ 乘出更多式，让“方程数追上单项式数”，把一个\emph{非线性消元}问题翻译成“方阵何时奇异”这一\emph{线性代数}问题（$\det=0$）。代价是引入了高次单项式（$x_2^3$ 等）、矩阵变大，且 $\det A$ 的次数可能\emph{高于}真实解数（出现增根，须回代验证）。这套“乘单项式 $\to$ 配平 $\to$ 取行列式”正是下面通用 6R 逆解的引擎，只是 6R 情形要嵌套用两次、矩阵到 $12\times12$。
\end{insight}

\paragraph{通用 6R 的消元流程（八步骨架）。}
把上法搬到 6R 臂。运动学先用 Denavit--Hartenberg 参数（\cref{def:ak-dh}）写成 $g_{st}(\boldsymbol{\theta})=g_{l_0l_1}(\theta_1)\cdots g_{l_5l_6}(\theta_6)=g_d$，每个 $g_{l_{i-1}l_i}$ 取 \cref{eq:ak-dh-matrix} 形（含 $c_i=\cos\theta_i,s_i=\sin\theta_i$）。把 $g_{l_5l_6}^{-1},g_{l_0l_1}^{-1},g_{l_1l_2}^{-1}$ 移到右端，得\emph{巧妙重排}\cite{raghavan1993inverse,murray1994mathematical}
\begin{equation}
g_{l_2l_3}(\theta_3)g_{l_3l_4}(\theta_4)g_{l_4l_5}(\theta_5)=g_{l_1l_2}^{-1}(\theta_2)g_{l_0l_1}^{-1}(\theta_1)\,g_d\,g_{l_6t}^{-1}g_{l_5l_6}^{-1}(\theta_6).
\label{eq:ak-6r-rearrange}
\end{equation}
之所以如此排，是为下面消元铺路。各步\emph{逻辑}如下（中间系数定义纯属机械记账，详见所引论文）：
\begin{description}
  \item[步 1--2] 选 DH 参数使 \cref{eq:ak-6r-rearrange} 右端第 3、4 列\emph{不依赖} $\theta_6$；比对两端第 3、4 列得 $6$ 个关于 $\theta_1,\dots,\theta_5$ 的方程 EQ1--EQ6，每个在 $s_i,c_i$ 中至多\emph{一次}（unary）。
  \item[步 3--4] 把 EQ1--EQ6 整理成两组三式（含向量 $\mathbf{p},\mathbf{l},\mathbf{n},\mathbf{h},\mathbf{f},\mathbf{r}$ 的双线性式）；再由 $\mathbf{l},\mathbf{p},\mathbf{p}\!\cdot\!\mathbf{p},\mathbf{p}\!\cdot\!\mathbf{l},\mathbf{p}\!\times\!\mathbf{l},(\mathbf{p}\!\cdot\!\mathbf{p})\mathbf{l}-2(\mathbf{p}\!\cdot\!\mathbf{l})\mathbf{p}$ 这 $3+3+1+1+3+3=14$ 个同型表达式，合成
  \begin{equation}
  P(s_3,c_3)\,\mathbf{x}_{45}=Q\,\mathbf{x}_{12},\quad P\in\mathbb{R}^{14\times9},\ Q\in\mathbb{R}^{14\times8}\ \text{常阵},
  \label{eq:ak-6r-PQ}
  \end{equation}
  其中 $\mathbf{x}_{45}=[s_4s_5,s_4c_5,c_4s_5,c_4c_5,s_4,c_4,s_5,c_5,1]^\top$、$\mathbf{x}_{12}=[s_1s_2,\dots,c_2]^\top$。
  \item[步 5] 设 $\mathrm{rank}\,Q=8$，用 14 式中任取 8 式解出 $\mathbf{x}_{12}$ 的 8 个分量（用 $\mathbf{x}_{45}$ 表示），回代余下 6 式得 $\Sigma(c_3,s_3)\,\mathbf{x}_{45}=\mathbf{0}$，$\Sigma\in\mathbb{R}^{6\times9}$。
  \item[步 6] 用正切半角代换 $s_i=\tfrac{2x_i}{1+x_i^2},c_i=\tfrac{1-x_i^2}{1+x_i^2}$（$x_i=\tan\tfrac{\theta_i}{2}$，$i=3,4,5$）清分母，把 $\Sigma$ 化为关于 $x_3$ 多项式系数的 $\Sigma_1(x_3)$。
  \item[步 7] 第二次 dialytic 消元：把 6 式乘 $x_4$ 再并入原 6 式，消去 $x_4,x_5$，配平成 $12\times12$ 方阵 $\bar{\Sigma}(x_3)$（如 \cref{eq:ak-dialytic-A} 之放大版）。
  \item[步 8] 解 $\det\bar{\Sigma}(x_3)=0$ 得 $x_3$（\emph{即 $\theta_3$}），回代零空间得 $x_4,x_5$，再由 \cref{eq:ak-6r-PQ} 得 $\theta_1,\theta_2$；最后回到原运动学，取任一点 $\mathbf{p}$ 写 $e^{[\mathcal{S}_6]\theta_6}\mathbf{p}=e^{-[\mathcal{S}_5]\theta_5}\cdots e^{-[\mathcal{S}_1]\theta_1}g_d\mathbf{M}^{-1}\mathbf{p}$，用子问题 1（\cref{thm:ak-pk-sub1}）解 $\theta_6$。
\end{description}

\begin{theorem}[通用 6R 逆解的尖锐上界 16（Lee--Liang）]\label{thm:ak-6r-bound}
通用六自由度串联（6R）机械臂的逆运动学，经上述两重 dialytic 消元归结为关于 $x_3=\tan(\theta_3/2)$ 的\textbf{16 次}多项式 $\det\bar{\Sigma}(x_3)=0$；从而\emph{至多有 16 组}解。该上界\textbf{尖锐}（存在确能达到 16 组实解的臂），由 Lee 与 Liang 首先证得，\emph{远优于}对原方程直接套 Bézout 定理所得的（松弛得多的）上界\cite{lee1988displacement,manocha1994realtime}。
\end{theorem}
\begin{proof}[说明]
$\det\bar{\Sigma}(x_3)$ 的次数由 $\bar{\Sigma}(x_3)\in\mathbb{R}^{12\times12}$ 各元关于 $x_3$ 的次数累加而得，化简后恰为 $16$（详细次数核算见所引论文，属机械计数）。$16$ 个根（含复根）对应 $\theta_3$ 的至多 $16$ 个值，每个回代唯一定出其余角，故逆解至多 16 组。\emph{尖锐性}＝构造出达到 16 实解的具体臂（Lee--Liang）。Bézout 界来自把三个三次多项式相乘的“盲目”上界，远大于 16；Lee--Liang 的功绩正是用 \cref{eq:ak-6r-rearrange} 的巧排把真实界压到 16。
\end{proof}

\begin{insight}[8 与 16：相交腕 vs 偏置腕的解数分水岭]\label{ins:ak-8-vs-16}
\cref{exam:ak-elbow-ik} 的肘式臂\emph{至多 8 解}，本节通用 6R \emph{至多 16 解}——差一倍，根源在\emph{结构}：肘式臂腕三轴\emph{交于一点}，使逆解能\emph{解耦}（先定腕心位置、再定姿态），每段几何子问题各贡献因子 2（$2\times2\times2=8$）；通用/偏置腕臂\emph{没有}这个交点，位置与姿态\emph{纠缠}，必须整体代数消元，多项式次数翻到 16。\textbf{这是“为可解性而设计”（\cref{thm:ak-pieper} 后的洞见）的量化代价}：球腕换来 8 解的闭式可解与数值便利；偏置腕（避奇异、增工作空间）换来更复杂的 16 解半解析求解。工程上：标准工业臂几乎都选球腕（8 解、实时闭式）；新型协作臂/避奇异设计才用偏置腕，配本节消元法或其数值化（广义特征值问题）实时求解\cite{manocha1994realtime}。
\end{insight}

\begin{note}[消元法的工程软肋]\label{note:ak-elim-caveat}
本法虽通用，却有实在的数值短板\cite{manocha1994realtime}：$\det\bar{\Sigma}(x_3)$ 涉及\emph{病态}矩阵的行列式/根，直接展开多项式再求根对系数扰动极敏感；故实用实现\emph{不}显式求多项式，而把 dialytic 消元转成\textbf{广义特征值问题} $\bar{\Sigma}(x_3)\mathbf{w}=\mathbf{0}$（$x_3$ 为特征值），借成熟的特征值算法改善条件数与速度——这对“在杂乱环境中实时算 IK 避障”至关重要。此外消元可能引入\emph{增根}（$\det$ 次数偶尔高于真实解数），须把每个根回代原运动学\emph{验证}。\textbf{要点}：通用 6R 有闭式（代数）解、至多 16 组，但稳健求解是数值工程，远不如球腕的子问题法干净——这也是工业界坚持球腕的现实理由。
\end{note}

\begin{practice}
\begin{enumerate}
  \item 对 \cref{exam:ak-dialytic}，验证 $A(x_1)$ 的上三行（原 $f_i$）在单项式 $x_2^3,x_2x_3$ 列为零、下三行（$x_2f_i$）在 $x_3,1$ 列为零，从而 $A$ 是“箭头/分块”结构；说明这一稀疏性从何而来。
  \item 为什么 dialytic 消元后 $\det A(x_1)$ 的次数\emph{可能高于}真实公共根数？举一个会出增根的情形，并说明实践中如何剔除（回代验证）。
  \item 通用 6R 的 16 解上界与肘式 8 解上界相差一倍。从“位置--姿态是否解耦”解释这个 2 倍，并说明它与 \cref{thm:ak-pieper} 球腕设计的关系。
  \item 为什么实用实现把 $\det\bar{\Sigma}(x_3)=0$ 转成广义特征值问题而非显式求多项式根？（提示：\cref{note:ak-elim-caveat} 的病态与条件数。）
\end{enumerate}
\end{practice}

% ════════════════════════════════════════════════════════════════
\section{速度运动学：雅可比}
\label{sec:ak-jacobian}

\paragraph{动机：从“位置”到“速度”，控制与力都在此层。}
正逆运动学处理\emph{位形}（位姿）。但控制要跟踪\emph{运动}（速度、加速度），力控要处理\emph{力}——这些都发生在\textbf{速度层}：末端速度 $\dot{\mathbf{x}}$ 与关节速度 $\dot{\mathbf{q}}$ 的关系。这个关系是线性的，其系数矩阵就是\textbf{雅可比} $\mathbf{J}$。雅可比是后续几乎所有内容（奇异、冗余、操作空间控制、力控、阻抗）的枢纽。

\begin{definition}[几何雅可比]\label{def:ak-jacobian}
末端空间速度（twist）$\boldsymbol{\mathcal{V}}=[\boldsymbol{\omega}_e;\,\mathbf{v}_e]\in\mathbb{R}^6$ 与关节速度 $\dot{\mathbf{q}}$ 的线性关系
\begin{equation}
  \boldsymbol{\mathcal{V}}=\mathbf{J}(\mathbf{q})\,\dot{\mathbf{q}},\qquad \mathbf{J}(\mathbf{q})\in\mathbb{R}^{6\times n}
  \label{eq:ak-jac-def}
\end{equation}
定义了\textbf{几何雅可比} $\mathbf{J}$，其第 $i$ 列 $\mathbf{J}_i$ 是“仅关节 $i$ 以单位速度运动、其余锁定”时末端的空间速度。
\end{definition}

\begin{theorem}[几何雅可比的逐列构造]\label{thm:ak-jac-columns}
几何雅可比的第 $i$ 列由关节 $i$ 的类型给出\cite{lynch2017modern,spong2006robot}（$\mathbf{z}_{i-1}$ 为关节 $i$ 轴在当前构型下的单位向量、$\mathbf{p}_{i-1}$ 为轴上一点（如 $O_{i-1}$）、$\mathbf{p}_e$ 为末端位置，均在 $\{0\}$ 下）：
\begin{equation}
  \boxed{\mathbf{J}_i=
  \begin{cases}
    \begin{bmatrix}\mathbf{z}_{i-1}\\[2pt]\mathbf{z}_{i-1}\times(\mathbf{p}_e-\mathbf{p}_{i-1})\end{bmatrix}, & \text{转动关节},\\[14pt]
    \begin{bmatrix}\mathbf{0}\\[2pt]\mathbf{z}_{i-1}\end{bmatrix}, & \text{移动关节}.
  \end{cases}}
  \label{eq:ak-jac-columns}
\end{equation}
（上 3 行为角速度部分、下 3 行为线速度部分。）
\end{theorem}
\begin{proof}
转动关节 $i$ 以单位角速度绕轴 $\mathbf{z}_{i-1}$ 转，给末端贡献角速度 $\boldsymbol{\omega}_e=\mathbf{z}_{i-1}$；末端作为该转动的刚体上一点，其线速度由刚体运动学（\cref{eq:poisson} 的 $\mathbf{v}=\boldsymbol{\omega}\times\mathbf{r}$）给出 $\mathbf{v}_e=\mathbf{z}_{i-1}\times(\mathbf{p}_e-\mathbf{p}_{i-1})$（$\mathbf{p}_e-\mathbf{p}_{i-1}$ 是从轴上点到末端的矢径）。移动关节 $i$ 以单位速度沿 $\mathbf{z}_{i-1}$ 平移，整段下游刚体平移、不转，故 $\boldsymbol{\omega}_e=\mathbf{0}$、$\mathbf{v}_e=\mathbf{z}_{i-1}$。叠加全部关节即 \cref{eq:ak-jac-def}。注意 $\mathbf{z}_{i-1},\mathbf{p}_{i-1},\mathbf{p}_e$ 均随构型 $\mathbf{q}$ 变，故 $\mathbf{J}=\mathbf{J}(\mathbf{q})$。
\end{proof}

\begin{note}[空间雅可比、物体雅可比与几何雅可比]\label{note:ak-jac-variants}
雅可比因“twist 表达在哪个系”“是否对应 POE 旋量”而有数个变体，本质同源\cite{lynch2017modern}：
\begin{itemize}
  \item \textbf{几何雅可比} $\mathbf{J}$（\cref{def:ak-jacobian}）：twist 各分量在 $\{0\}$ 下，$\mathbf{p}_e$ 取真实末端点，最直观，控制常用。
  \item \textbf{空间雅可比} $\mathbf{J}_s$：第 $i$ 列是\emph{当前构型下}的空间旋量 $\mathrm{Ad}_{(e^{[\mathcal{S}_1]\theta_1}\cdots e^{[\mathcal{S}_{i-1}]\theta_{i-1}})}\mathcal{S}_i$，由 POE（\cref{eq:ak-poe-space}）对 $\theta_i$ 求导直接得；其 twist 是物体相对 $\{0\}$ 的空间速度。
  \item \textbf{物体雅可比} $\mathbf{J}_b$：第 $i$ 列由物体形式 POE（\cref{eq:ak-poe-body}）给出，twist 在末端系 $\{n\}$ 下；与 $\mathbf{J}_s$ 经伴随相联 $\mathbf{J}_s=\mathrm{Ad}_{\mathbf{T}_{0n}}\mathbf{J}_b$（\cref{eq:Ad}）。
\end{itemize}
三者的角速度/线速度分量约定与参考点略有差别，互相用伴随 $\mathrm{Ad}_{\mathbf{T}}$ 精确换算（\cref{thm:se3-kinematics} 的 $\boldsymbol{\varpi}=\boldsymbol{\mathcal{J}}\dot{\boldsymbol{\xi}}$ 即李代数形式）。\textbf{工程陷阱}：库函数（Pinocchio、KDL）默认返回哪种雅可比、参考系是 \texttt{LOCAL} / \texttt{WORLD} / \texttt{LOCAL\_WORLD\_ALIGNED}，必须查清，否则力控/OSC 的 $\mathbf{J}^\top\mathbf{F}$ 会用错系（\cref{sec:ak-practice}）。
\end{note}

\begin{theorem}[解析雅可比]\label{thm:ak-jac-analytic}
若末端姿态用某\emph{最小参数} $\boldsymbol{\phi}$（如欧拉角）表示、任务向量 $\mathbf{x}=[\mathbf{p}_e;\boldsymbol{\phi}]$，则\textbf{解析雅可比} $\mathbf{J}_a$ 定义为 $\dot{\mathbf{x}}=\mathbf{J}_a\dot{\mathbf{q}}$，与几何雅可比相差一个仅依赖姿态表示的变换 $\mathbf{T}_\phi$\cite{spong2006robot}：
\begin{equation}
  \mathbf{J}_a=\begin{bmatrix}\mathbf{I}_3 & \mathbf{0}\\ \mathbf{0} & \mathbf{T}_\phi(\boldsymbol{\phi})^{-1}\end{bmatrix}\mathbf{J}',
  \label{eq:ak-jac-analytic}
\end{equation}
其中 $\mathbf{J}'$ 为线速度在上、角速度在下排列的几何雅可比，$\mathbf{T}_\phi$ 把 $\dot{\boldsymbol{\phi}}$ 映到角速度 $\boldsymbol{\omega}_e=\mathbf{T}_\phi\dot{\boldsymbol{\phi}}$。$\mathbf{T}_\phi$ 在姿态表示奇异处（如欧拉角万向锁）退化——这是\emph{表示奇异}，区别于真实的\emph{运动学奇异}。
\end{theorem}

\begin{theorem}[静力学对偶：$\boldsymbol{\tau}=\mathbf{J}^\top\mathbf{F}$]\label{thm:ak-static-duality}
机械臂\emph{静止}（或准静态）时，为在末端产生广义力（wrench）$\mathbf{F}=[\mathbf{m};\mathbf{f}]\in\mathbb{R}^6$（力矩在前、力在后，与 twist 配对），所需关节力矩为\cite{lynch2017modern}：
\begin{equation}
  \boxed{\boldsymbol{\tau}=\mathbf{J}^\top\mathbf{F}.}
  \label{eq:ak-static-duality}
\end{equation}
\end{theorem}
\begin{proof}
由\emph{虚功原理}：任一虚位移下，关节空间虚功等于任务空间虚功。关节虚位移 $\delta\mathbf{q}$ 经 \cref{eq:ak-jac-def} 给末端虚位移（twist 积分）$\delta\mathbf{x}=\mathbf{J}\delta\mathbf{q}$。虚功守恒 $\boldsymbol{\tau}^\top\delta\mathbf{q}=\mathbf{F}^\top\delta\mathbf{x}=\mathbf{F}^\top\mathbf{J}\delta\mathbf{q}$ 对任意 $\delta\mathbf{q}$ 成立，故 $\boldsymbol{\tau}^\top=\mathbf{F}^\top\mathbf{J}$，转置即 \cref{eq:ak-static-duality}。
\end{proof}

\begin{insight}[一个矩阵，两条对偶链路]
雅可比同时主管\emph{运动}与\emph{力}两条对偶链路：速度上 $\dot{\mathbf{x}}=\mathbf{J}\dot{\mathbf{q}}$（关节速 $\to$ 末端速），力上 $\boldsymbol{\tau}=\mathbf{J}^\top\mathbf{F}$（末端力 $\to$ 关节力矩）——\emph{方向相反、用同一个 $\mathbf{J}$ 的转置}。这一对偶是机械臂控制的支柱：\cref{ch:arm-ctrl} 的操作空间控制（$\boldsymbol{\tau}=\mathbf{J}^\top\mathbf{F}$）、阻抗控制（$\boldsymbol{\tau}=\mathbf{J}^\top(-\mathbf{K}\Delta\mathbf{x}-\mathbf{D}\dot{\mathbf{x}})$）、力位混合控制全靠它把“任务空间想要的力”翻译成“关节要发的力矩”。\cref{ch:arm-prac} 的 \texttt{arm\_dm} 力矩臂正是用 $\boldsymbol{\tau}=\mathbf{J}^\top\mathbf{F}$ 实现\emph{无需力传感器的真阻抗}（测位置 $\to$ 算 $\Delta\mathbf{x}$ $\to$ 经 $\mathbf{J}^\top$ 出力矩），与 P7 的 $\boldsymbol{\tau}=\mathbf{J}^\top\mathbf{F}$（\cref{sec:qk-jacobian}，足端力 $\to$ 关节力矩）同源。
\end{insight}

\begin{practice}
\begin{enumerate}
  \item 对 2R 平面臂求几何雅可比（$2\times2$ 的线速度块即可），验证 $\mathbf{J}=\big[\begin{smallmatrix}-a_1 s_1-a_2 s_{12} & -a_2 s_{12}\\ a_1 c_1+a_2 c_{12} & a_2 c_{12}\end{smallmatrix}\big]$（$s_{12}=\sin(\theta_1+\theta_2)$ 等）。
  \item 用 \cref{thm:ak-static-duality} 算：上题 2R 臂末端受竖直向下力 $f$ 时，两关节各需多大力矩平衡？验证它等于“重力矩”的直觉。
  \item 解释为何解析雅可比 $\mathbf{J}_a$ 会在欧拉角万向锁处奇异，而几何雅可比 $\mathbf{J}$ 在同一构型可能良好——这两种“奇异”有何本质不同？
\end{enumerate}
\end{practice}

% ════════════════════════════════════════════════════════════════
\section{奇异性与可操作度}
\label{sec:ak-singularity}

\paragraph{动机：有些构型下，臂“使不上劲”。}
雅可比 $\mathbf{J}(\mathbf{q})$ 随构型变。在某些构型，$\mathbf{J}$ \emph{掉秩}——末端在某个方向上\textbf{无论关节怎么动都动不了}（瞬时不可达），或反过来要某个末端速度需\emph{无穷大}的关节速度。这叫\textbf{运动学奇异}（singularity）。奇异是机械臂的“暗礁”：控制律在此失效、IK 数值法发散、力在某方向无法施加。本节讲如何\emph{判别}、\emph{量化}、\emph{绕过}奇异。

\begin{definition}[运动学奇异与可操作度]\label{def:ak-singularity}
构型 $\mathbf{q}$ 是\textbf{奇异}的，当 $\mathbf{J}(\mathbf{q})$ 不满秩（$\mathrm{rank}\,\mathbf{J}<\min(6,n)$）。对非冗余臂（$n=6$），等价于 $\det\mathbf{J}=0$。\textbf{可操作度}（manipulability measure，Yoshikawa）量化“离奇异多远”\cite{lynch2017modern}：
\begin{equation}
  \boxed{w(\mathbf{q})=\sqrt{\det\!\big(\mathbf{J}\mathbf{J}^\top\big)}=\sigma_1\sigma_2\cdots\sigma_m\ge0,}
  \label{eq:ak-manip}
\end{equation}
$\sigma_k$ 为 $\mathbf{J}$ 的奇异值。$w=0\iff$ 奇异；$w$ 越大越“灵活”。
\end{definition}

\begin{insight}[可操作度椭球：末端能动多快的“方向地图”]
$w$ 的几何意义由\textbf{可操作度椭球}给出\cite{lynch2017modern}：单位关节速度球 $\|\dot{\mathbf{q}}\|=1$ 经 $\dot{\mathbf{x}}=\mathbf{J}\dot{\mathbf{q}}$ 映成任务空间一个椭球，其主轴方向是 $\mathbf{J}\mathbf{J}^\top$ 的特征向量、半轴长是奇异值 $\sigma_k$。椭球\emph{越圆}越各向同性（各方向同样灵活）、\emph{越扁}越接近奇异（最短轴方向几乎动不了）。$w=\sqrt{\det(\mathbf{J}\mathbf{J}^\top)}$ 正比于椭球\emph{体积}。\textbf{这给了规划一个目标}：让臂工作在 $w$ 大的“灵巧区”，远离 $w\to0$ 的奇异——\cref{ch:arm-traj} 可把 $w$ 作为轨迹优化的代价项，\cref{sec:ak-redundancy} 可用零空间最大化 $w$。
\end{insight}

\begin{figure}[ht]
\centering
\begin{tikzpicture}[scale=1.0,>=Stealth]
% good config: round-ish ellipse
\begin{scope}[shift={(0,0)}]
  \draw[thick, main!70!black, fill=main!8] (0,0) ellipse (0.95 and 0.7);
  \draw[->,thick] (0,0)--(0.95,0); \draw[->,thick] (0,0)--(0,0.7);
  \node at (0,-1.15) {\scriptsize 良态：$w$ 大（椭球近圆）};
\end{scope}
% near-singular: flat ellipse
\begin{scope}[shift={(4,0)}]
  \draw[thick, second!70!black, fill=second!8] (0,0) ellipse (1.15 and 0.16);
  \draw[->,thick] (0,0)--(1.15,0); \draw[->,thick] (0,0)--(0,0.16);
  \node at (0,-1.15) {\scriptsize 近奇异：$w\to0$（椭球塌扁）};
  \node[second!70!black] at (0,0.5) {\scriptsize 短轴方向几乎不可动};
\end{scope}
\end{tikzpicture}
\caption{可操作度椭球。良态构型椭球近圆（各向灵活）；近奇异构型椭球沿某方向塌扁，$w=\sqrt{\det(\mathbf{J}\mathbf{J}^\top)}\to0$，该方向上末端瞬时不可动。}
\label{fig:ak-manip-ellipsoid}
\end{figure}

\begin{insight}[趋近奇异 = 椭球某轴连续塌为线 = 那个方向被“冻结”]
把 \cref{fig:ak-manip-ellipsoid} 看成一段\emph{连续过程}而非两张快照：当构型连续逼近奇异，可操作度椭球的\emph{最短半轴}（长度恰为最小奇异值 $\sigma_{\min}$）\textbf{连续缩短}，在奇异处\emph{塌缩为零}——椭球退化成一张薄片乃至一条线。几何直觉一句话：\textbf{那条塌掉的轴所指方向，正是末端“无论关节怎么转都瞬时动不了”的方向}（沿它末端速度的可达上界 $\to0$）；反过来要在该方向得到哪怕一点点末端速度，所需关节速度 $\propto1/\sigma_{\min}\to\infty$，这正是 \cref{thm:ak-dls} 里 $\mathbf{J}^{+}$ “爆炸”的椭球图像。于是 $\det\mathbf{J}=0$（代数）、椭球塌一根轴（几何）、某方向瞬时不可达（物理）是\emph{同一件事}的三种说法（Yoshikawa）\cite{yoshikawa1985manipulability}——这也解释了为何用\emph{体积} $w=\prod_k\sigma_k$ 而非长度做指标：只要\emph{任何一根}轴塌，体积即归零，奇异无所遁形。
\end{insight}

\paragraph{典型奇异构型。}
对带球腕 6R 臂，奇异分三类\cite{spong2006robot}：(1) \textbf{腕奇异}（$q_5=0$，腕两轴 $z_4,z_6$ 共线，损失一个姿态自由度，最常遇）；(2) \textbf{肘奇异}（臂\emph{完全伸直}或完全折叠，肘失去径向运动能力）；(3) \textbf{肩奇异}（腕心落在第一关节轴 $z_0$ 上，绕基座 yaw 不改变腕心位置）。这些都是 $\det\mathbf{J}=0$ 的具体几何。

\begin{theorem}[阻尼最小二乘（DLS）伪逆]\label{thm:ak-dls}
在 $\dot{\mathbf{x}}=\mathbf{J}\dot{\mathbf{q}}$ 上反解关节速度时，奇异附近 $\mathbf{J}^{+}=\mathbf{J}^\top(\mathbf{J}\mathbf{J}^\top)^{-1}$ 因 $\mathbf{J}\mathbf{J}^\top$ 病态而\emph{爆炸}。\textbf{阻尼最小二乘}（damped least squares，Levenberg--Marquardt 正则）用 $\lambda>0$ 牺牲少量精度换数值稳定\cite{spong2006robot,lynch2017modern}：
\begin{equation}
  \boxed{\dot{\mathbf{q}}=\mathbf{J}^\top\big(\mathbf{J}\mathbf{J}^\top+\lambda^2\mathbf{I}\big)^{-1}\dot{\mathbf{x}}.}
  \label{eq:ak-dls}
\end{equation}
它等价于最小化 $\|\mathbf{J}\dot{\mathbf{q}}-\dot{\mathbf{x}}\|^2+\lambda^2\|\dot{\mathbf{q}}\|^2$，即在“贴合任务速度”与“关节速度不过大”之间折中。
\end{theorem}
\begin{proof}
最小化 $g(\dot{\mathbf{q}})=\tfrac12\|\mathbf{J}\dot{\mathbf{q}}-\dot{\mathbf{x}}\|^2+\tfrac{\lambda^2}{2}\|\dot{\mathbf{q}}\|^2$，令梯度 $\mathbf{J}^\top(\mathbf{J}\dot{\mathbf{q}}-\dot{\mathbf{x}})+\lambda^2\dot{\mathbf{q}}=\mathbf{0}$，得 $(\mathbf{J}^\top\mathbf{J}+\lambda^2\mathbf{I})\dot{\mathbf{q}}=\mathbf{J}^\top\dot{\mathbf{x}}$，即 $\dot{\mathbf{q}}=(\mathbf{J}^\top\mathbf{J}+\lambda^2\mathbf{I})^{-1}\mathbf{J}^\top\dot{\mathbf{x}}$；用矩阵恒等式 $(\mathbf{J}^\top\mathbf{J}+\lambda^2\mathbf{I})^{-1}\mathbf{J}^\top=\mathbf{J}^\top(\mathbf{J}\mathbf{J}^\top+\lambda^2\mathbf{I})^{-1}$（push-through，对任意 $\lambda>0$ 两侧均良定）即得 \cref{eq:ak-dls}。从奇异值看：$\sigma_k$ 经 DLS 变为 $\sigma_k/(\sigma_k^2+\lambda^2)$，在 $\sigma_k\to0$ 时该因子 $\to\sigma_k/\lambda^2\to0$（不再爆炸），而 $\sigma_k\gg\lambda$ 时 $\approx1/\sigma_k$（几乎不损精度）——$\lambda$ 正是“在多小的奇异值上开始阻尼”的阈值。
\end{proof}

\begin{note}[复用 P7 已证的伪逆/DLS：本节不重证]\label{note:ak-dls-reuse}
MP 伪逆 $\mathbf{J}^{+}=\mathbf{J}^\top(\mathbf{J}\mathbf{J}^\top)^{-1}$ 的“最小范数”性质、零空间投影、阻尼最小二乘的稳健化，已在 \cref{sec:qk-nullspace}（四足部）用拉格朗日乘子法\emph{完整证明}（\cref{thm:qk-mp-minnorm,thm:qk-nullspace-proj}）。本节\textbf{直接复用其结论}，不重复证明，只补“机械臂奇异分类 + 可操作度椭球”这一专门视角。读者需要伪逆/DLS 的严格推导时请回看 \cref{sec:qk-nullspace}。
\end{note}

\begin{pitfall}[陷阱：在力矩级硬刚奇异 vs 在速度级稳健化]
DLS（\cref{eq:ak-dls}）工作在\emph{速度/运动学层}（无质量矩阵），在奇异附近稳健、可靠。一个深刻的工程教训来自 \cref{ch:arm-prac}：试图把同类“正则化”搬到\emph{力矩级}的操作空间控制（OSC，给 $\mathbf{J}\mathbf{M}^{-1}\mathbf{J}^\top$ 加阻尼）时，正则化会\textbf{破坏动力学一致投影的投影性}（特征值不再是 $\{0,1\}$），反而让任务空间漏入扰动。结论：\textbf{奇异/冗余的稳健化在速度级做最干净}；力矩级 OSC 对质量矩阵条件数极敏感，须谨慎（\cref{ch:arm-dyn} 的条件数分析、\cref{ch:arm-ctrl} 的 OSC 完整推导）。这是“同一个正则化思想，在不同层效果迥异”的典型。
\end{pitfall}

\paragraph{案例：折叠臂的子雅可比奇异。}
\cref{ch:arm-prac} 的力位混合控制实践中，遇到一个具体奇异陷阱：在臂\emph{折叠}的构型下，若按“位控用某几列、力控用另几列”切分雅可比，切出的\emph{子雅可比}恰好奇异（折叠位下相关关节轴近共线），导致控制飞车。诊断历程揭示的正解是：\textbf{用完整 6 列雅可比做笛卡尔控制、让冗余腕安全进入零空间}，而非按关节切分子雅可比（详见 \cref{ch:arm-ctrl} 力位混合、\cref{ch:arm-prac} 诊断方法论）。这印证了 \cref{def:ak-singularity}：奇异是\emph{构型相关}的，规整构型下良态的臂在折叠/伸直处会突然奇异。

\begin{practice}
\begin{enumerate}
  \item 对 2R 平面臂，算 $\det\mathbf{J}$，证明奇异恰发生在 $\theta_2=0$ 或 $\pi$（臂伸直或完全折叠），并说明此时末端在哪个方向不可动。
  \item 取 $\lambda=0.1$，画出 DLS 因子 $\sigma/(\sigma^2+\lambda^2)$ 随 $\sigma$ 的曲线，标出“几乎不损精度”区（$\sigma\gg\lambda$）与“被阻尼压制”区（$\sigma\lesssim\lambda$）。
  \item 解释为何 $w=\prod_k\sigma_k$ 而非 $\min_k\sigma_k$ 作为可操作度——前者对“某一个 $\sigma$ 变小”的敏感性如何？讨论两者作为奇异指标的优劣。
\end{enumerate}
\end{practice}

% ════════════════════════════════════════════════════════════════
\section{冗余与零空间解析}
\label{sec:ak-redundancy}

\paragraph{动机：自由度“富余”是负担还是资源？}
当关节数 $n$ 多于任务维 $m$（$n>m$），臂\textbf{冗余}：达成同一末端任务有\emph{无穷多}关节构型。看似累赘，实为资源——多出的 $n-m$ 维自由度可在\emph{不影响主任务}的前提下做次要任务：避障、避关节限位、避奇异（最大化 §1.6 的 $w$）、优化能耗。本节讲如何用\textbf{零空间}系统地调度这些富余自由度。

\paragraph{反面：忽略冗余＝浪费 + 不可控的内部运动。}
若对冗余臂仍用某个固定右逆求 IK，多出的自由度被\emph{隐式}固定，机器人可能在不知不觉中漂向关节限位或奇异，也白白浪费了避障能力。正确做法是\emph{显式}用零空间投影把次要任务编码进去。

\begin{definition}[降维造冗余]\label{def:ak-redundancy}
即便\emph{非冗余}臂（$n=6$），若主任务只约束\emph{部分}维度（如只控 3D 位置、不控姿态），任务维 $m<n$，则系统对该任务\emph{冗余}，零空间维数 $n-m>0$。例：6 DoF 臂只控末端 3D 位置（$m=3$），留 $6-3=3$ 维零空间做\emph{自运动}（self-motion，关节动而末端位置不动）。
\end{definition}

\begin{theorem}[resolved-rate 冗余解析]\label{thm:ak-resolved-rate}
冗余下关节速度的通解为“主任务特解 + 零空间齐次解”\cite{lynch2017modern}：
\begin{equation}
  \boxed{\dot{\mathbf{q}}=\underbrace{\mathbf{J}^{+}\big(\dot{\mathbf{x}}_d+\mathbf{K}\mathbf{e}\big)}_{\text{主任务（闭环）}}+\underbrace{\mathbf{N}\,\dot{\mathbf{q}}_{\text{sec}}}_{\text{零空间次任务}},\qquad \mathbf{N}=\mathbf{I}-\mathbf{J}^{+}\mathbf{J},}
  \label{eq:ak-resolved-rate}
\end{equation}
其中 $\mathbf{e}=\mathbf{x}_d-\mathbf{x}$ 为任务误差、$\mathbf{K}\succ0$ 反馈增益（使主任务渐近收敛、抑制数值漂移）、$\dot{\mathbf{q}}_{\text{sec}}$ 为次任务期望关节速度、$\mathbf{N}$ 把它投影到零空间。无论 $\dot{\mathbf{q}}_{\text{sec}}$ 取何值，主任务恒满足。
\end{theorem}
\begin{proof}
左乘 $\mathbf{J}$：$\mathbf{J}\dot{\mathbf{q}}=\mathbf{J}\mathbf{J}^{+}(\dot{\mathbf{x}}_d+\mathbf{K}\mathbf{e})+\mathbf{J}\mathbf{N}\dot{\mathbf{q}}_{\text{sec}}$。由 \cref{sec:qk-nullspace} 已证 $\mathbf{J}\mathbf{J}^{+}=\mathbf{I}$、$\mathbf{J}\mathbf{N}=\mathbf{0}$（\cref{thm:qk-nullspace-proj}），得 $\dot{\mathbf{x}}=\dot{\mathbf{x}}_d+\mathbf{K}\mathbf{e}$，即 $\dot{\mathbf{e}}+\mathbf{K}\mathbf{e}=\mathbf{0}$（用 $\dot{\mathbf{e}}=\dot{\mathbf{x}}_d-\dot{\mathbf{x}}$），主任务误差指数收敛、与 $\dot{\mathbf{q}}_{\text{sec}}$ 无关。零空间项只在 $\mathbf{N}$ 的像（即 $\mathcal{N}(\mathbf{J})$）内运动，不污染主任务——这就是“免费内部运动”（\cref{ins:qk-nullspace-free}）在 resolved-rate 框架下的兑现。
\end{proof}

\begin{insight}[零空间＝“在完美做主任务的前提下，尽力做次任务”]
\cref{eq:ak-resolved-rate} 的哲学（\cref{sec:qk-nullspace} 已深入）：主任务\emph{严格}满足（占满 $\mathbf{J}$ 的行空间），次任务\emph{尽力}近似（被 $\mathbf{N}$ 投到零空间、取其零空间分量）。这是\textbf{严格任务优先级}（strict hierarchy）的几何实现，也是 \cref{sec:ctrl-robot-wbc}（全身控制 WBC）多任务分层 QP 的运动学雏形。常用次任务 $\dot{\mathbf{q}}_{\text{sec}}$ 取某标量势 $H(\mathbf{q})$ 的梯度上升 $\dot{\mathbf{q}}_{\text{sec}}=k\nabla H$：$H=w(\mathbf{q})$ 则避奇异、$H=-\sum(q_i-q_i^{\text{mid}})^2$ 则避限位、$H=$ 到障碍距离则避障（\cref{ch:arm-traj} 的反应式避障即用零空间或 CBF）。
\end{insight}

\paragraph{案例：运动学级零空间——EE 锁 0.00mm + J1 主导自运动。}
\cref{ch:arm-prac} 的一个高价值实验把 \cref{eq:ak-resolved-rate} 落到 \texttt{arm\_dm} 臂：主任务降维到 3D EE 位置（$m=3$），留 3 维零空间。结果\textbf{末端位置精确锁定 $\|\Delta\mathbf{x}\|=0.00\,$mm}（仅积分误差量级），同时零空间驱动姿态自运动——\emph{臂绕固定工具点重构姿态而工具不动}。更深的几何洞察：自运动总由 \textbf{$J_1$（基座 yaw）主导}（实测 $q_1$ 摆约 $0.44\,$rad $\approx25^\circ$ 而 EE 不动），因为该臂\emph{末端几乎落在基座 yaw 轴上}——对“3D 位置”这一主任务，绕基座 yaw 是最自由的零空间方向。这把“臂绕固定工具重构”从现象升为几何洞察。\textbf{关键对照}：此实验在\emph{运动学/速度级}（用 $\mathbf{N}=\mathbf{I}-\mathbf{J}^{+}\mathbf{J}$，无质量矩阵 $\mathbf{M}^{-1}$）做，故\emph{干净}；其\emph{力矩级}对应（动力学一致 OSC 零空间 $\mathbf{N}=\mathbf{I}-\mathbf{J}^\top(\mathbf{J}\mathbf{M}^{-1}\mathbf{J}^\top)^{-1}\mathbf{J}\mathbf{M}^{-1}$，Khatib）因该臂质量矩阵\emph{病态}而 EE 漂 60--110mm——这一层次对照（运动学级干净 vs 力矩级受病态限制）留到 \cref{ch:arm-dyn}（条件数）与 \cref{ch:arm-ctrl}（OSC 完整推导）深入。本节立住结论：\textbf{异质惯量臂的冗余解析正解＝运动学/速度级。}

\begin{practice}
\begin{enumerate}
  \item 对一个 3R 平面臂（$n=3$）只控末端 2D 位置（$m=2$），写出 $\mathbf{N}=\mathbf{I}-\mathbf{J}^{+}\mathbf{J}$（$3\times3$，秩 1），并说明零空间自运动的几何含义。
  \item 取次任务为“关节居中” $H=-\tfrac12\sum_i(q_i-q_i^{\text{mid}})^2$，写出 $\dot{\mathbf{q}}_{\text{sec}}=\nabla H$ 并代入 \cref{eq:ak-resolved-rate}，解释它如何在不动末端的前提下把关节推离限位。
  \item 为什么 \cref{eq:ak-resolved-rate} 主任务项要加反馈 $\mathbf{K}\mathbf{e}$ 而非只用前馈 $\mathbf{J}^{+}\dot{\mathbf{x}}_d$？（提示：数值积分漂移、$\mathbf{J}^{+}$ 的近似性。）
\end{enumerate}
\end{practice}

% ════════════════════════════════════════════════════════════════
\section{并联机构运动学}
\label{sec:ak-parallel}

\paragraph{动机：把“串联”反过来读。}
本章迄今讲的都是\emph{串联}机械臂——连杆一根接一根的开链（\cref{def:ak-serial}）。但工程中另有一大类机构：末端平台经\emph{多条}运动链\emph{并行}连到基座，构成\textbf{并联机构}（parallel manipulator）。Stewart--Gough 飞行模拟平台、Delta 高速分拣机器人、并联机床都属此类。Siciliano 在介绍机械臂结构时即指出\cite{siciliano2009robotics}：并联（闭链）几何相对开链的根本优势是\emph{结构刚度高、可达极高操作速度}（六条腿分担负载、电机可置于基座而非随动），代价是\emph{工作空间小、构型耦合强}。本节把串联那套（FK/IK/雅可比/奇异）\emph{对偶}地搬到并联上，会发现一处处“正难者并联易、并联难者串联易”的镜像。

\paragraph{反面：照搬串联的“连乘 FK”会卡死。}
串联 FK 是相邻位姿连乘（\cref{eq:ak-poschain}）、一遍算完。但并联平台的位姿\emph{不能}由各腿长简单连乘得到——每条腿都对平台施加一个约束，平台位姿须\emph{同时满足}全部约束方程，这是一组\emph{耦合非线性方程}。盲目套用串联思路会发现：并联的\textbf{正运动学反而是难的}。这正是并联与串联最醒目的对偶。

\begin{definition}[并联机构与两类典型]\label{def:ak-parallel}
\textbf{并联机构}是动平台（moving platform，末端）与定平台（base）之间由 $\ge2$ 条独立运动链（“腿”，limb）连接、构成一个或多个\emph{闭合运动链}的机构。两类工程典范\cite{siciliano2009robotics,merlet2006parallel}：
\begin{itemize}
  \item \textbf{Stewart--Gough 平台}（6\nobreakdash-UPS）：动平台经 6 条可变长腿连到基座，每腿为“万向节 U + 移动副 P（驱动）+ 球副 S”。6 条腿长 $\boldsymbol{\ell}=(\ell_1,\dots,\ell_6)$ 为驱动变量，平台具完整 6 自由度位姿。广用于飞行/驾驶模拟器、六维力台、精密定位。
  \item \textbf{Delta 机构}：动平台经 3 条含\emph{平行四边形}的腿连到基座，每条腿由一个基座转动副（驱动）带动一组平行四边形。平行四边形约束使动平台\emph{始终平动}（姿态恒定），保留 3 个平移自由度。结构轻、惯量小，是高速取放（pick-and-place）的主力（如包装、分拣）。
\end{itemize}
\end{definition}

\begin{insight}[串联--并联的运动学对偶：正逆难易恰好互换]\label{ins:ak-parallel-duality}
并联机构在运动学上几乎处处是串联的\emph{镜像}\cite{merlet2006parallel}：
\begin{center}\small
\begin{tabular}{lll}
\toprule
 & 串联开链 & 并联闭链 \\
\midrule
正运动学（驱动 $\to$ 末端） & \emph{易}：连乘一遍即得 & \emph{难}：耦合非线性、\textbf{多解}、常需数值 \\
逆运动学（末端 $\to$ 驱动） & \emph{难}：解三角方程组、多解 & \emph{易}：每条腿\textbf{独立}闭式反解 \\
工作空间 & 大 & 小（受各腿行程 + 干涉限制） \\
刚度 / 速度 & 较低 / 较低 & \textbf{高} / \textbf{高}（负载分担、电机置基座） \\
奇异的危险方向 & 失\emph{自由度}（动不了） & 多一类失\emph{约束}（\textbf{失控、撑不住}） \\
\bottomrule
\end{tabular}
\end{center}
直觉根源：串联是“一串自由度顺次叠加”，给定关节角末端唯一（正易）、但反求要解一串耦合三角方程（逆难）；并联是“多条约束并行交汇”，给定末端位姿每条腿\emph{各自}量长度即可（逆易），但反过来由腿长定平台位姿要解所有约束的交（正难）。
\end{insight}

\paragraph{逆运动学：并联的“送分题”。}
以 Stewart--Gough 为例。设动平台位姿 $\mathbf{T}=\big[\begin{smallmatrix}\mathbf{R}&\mathbf{p}\\\mathbf{0}^\top&1\end{smallmatrix}\big]\in\SE(3)$，第 $i$ 条腿在基座的铰点为 $\mathbf{b}_i$（基座系 $\{0\}$ 下常量）、在动平台的铰点为 $\mathbf{a}_i$（平台系下常量）。则该腿在 $\{0\}$ 下的矢量为 $\mathbf{p}+\mathbf{R}\mathbf{a}_i-\mathbf{b}_i$，腿长直接得\cite{merlet2006parallel}：
\begin{equation}
  \ell_i=\big\|\,\mathbf{p}+\mathbf{R}\,\mathbf{a}_i-\mathbf{b}_i\,\big\|,\qquad i=1,\dots,6.
  \label{eq:ak-stewart-ik}
\end{equation}
\textbf{六条腿彼此独立}、各取一次范数即得，无迭代、无多解歧义——这就是并联 IK“易”的全部内容。Delta 的逆解同理：给定动平台位置，每条腿的平行四边形几何给出一个\emph{平面}两连杆问题，闭式（余弦定理）解出该腿驱动转角，三腿独立。

\begin{insight}[正运动学才是并联的“拦路虎”]\label{ins:ak-parallel-fk}
反过来，已知 6 个腿长 $\boldsymbol{\ell}$ 求平台位姿 $\mathbf{T}$（并联\emph{正}运动学），要解 \cref{eq:ak-stewart-ik} 的 6 个方程联立——这是一组高度耦合的非线性方程，一般\emph{无闭式}、需数值迭代（牛顿法），且\textbf{解不唯一}：通用 6\nobreakdash-6 Stewart 平台的正运动学最多可有 \textbf{40 组}实解\cite{merlet2006parallel}（即同一组腿长对应至多 40 个平台位姿）。这与串联“正易逆难”形成完整对偶。工程上常加装额外传感（如被动腿编码器、惯导）以\emph{实时唯一确定}当前位姿，回避求解 40 解的难题。
\end{insight}

\paragraph{并联雅可比：由约束方程隐式给出。}
串联雅可比可逐列\emph{显式}写出（\cref{eq:ak-jac-columns}）；并联则不然——驱动变量与末端位姿被\emph{约束方程}隐式绑定，雅可比要\emph{两个矩阵}才说清。

\begin{theorem}[并联机构的隐式雅可比]\label{thm:ak-parallel-jac}
设驱动变量 $\mathbf{q}\in\mathbb{R}^n$（如 Stewart 的腿长 $\boldsymbol{\ell}$）、末端位姿用 $\mathbf{x}\in\mathbb{R}^m$（位置 + 姿态参数）描述，约束方程组（每条腿一式，如 \cref{eq:ak-stewart-ik}）写成隐式形式 $\mathbf{f}(\mathbf{x},\mathbf{q})=\mathbf{0}$。对时间求导得速度级关系\cite{merlet2006parallel}：
\begin{equation}
  \boxed{\;\mathbf{J}_x\,\dot{\mathbf{x}}=\mathbf{J}_q\,\dot{\mathbf{q}},\qquad
  \mathbf{J}_x=\frac{\partial\mathbf{f}}{\partial\mathbf{x}},\quad
  \mathbf{J}_q=-\frac{\partial\mathbf{f}}{\partial\mathbf{q}}.\;}
  \label{eq:ak-parallel-jac}
\end{equation}
当 $\mathbf{J}_q$ 可逆时，\emph{逆}雅可比（末端速 $\to$ 驱动速）为 $\mathbf{J}^{-1}=\mathbf{J}_q^{-1}\mathbf{J}_x$；当 $\mathbf{J}_x$ 可逆时，\emph{正}雅可比为 $\mathbf{J}=\mathbf{J}_x^{-1}\mathbf{J}_q$。整体雅可比 $\dot{\mathbf{q}}=\mathbf{J}_q^{-1}\mathbf{J}_x\,\dot{\mathbf{x}}$。
\end{theorem}
\begin{proof}
对 $\mathbf{f}(\mathbf{x},\mathbf{q})=\mathbf{0}$ 全微分（隐函数定理）：$\frac{\partial\mathbf{f}}{\partial\mathbf{x}}\dot{\mathbf{x}}+\frac{\partial\mathbf{f}}{\partial\mathbf{q}}\dot{\mathbf{q}}=\mathbf{0}$，移项即 \cref{eq:ak-parallel-jac}。对 Stewart 平台，把 \cref{eq:ak-stewart-ik} 平方 $\ell_i^2=\|\mathbf{p}+\mathbf{R}\mathbf{a}_i-\mathbf{b}_i\|^2$ 求导，可得 $\mathbf{J}_q=\mathrm{diag}(\ell_i)$（对角、几乎总可逆——故逆运动学雅可比简单），而 $\mathbf{J}_x$ 的第 $i$ 行是腿单位向量 $\hat{\mathbf{s}}_i$ 与其力臂构成的\emph{腿旋量}，$\mathbf{J}_x^{-1}$（正雅可比）才是难点——它在某些位形会奇异，对应下文的直接运动学奇异。
\end{proof}

\begin{insight}[“两个雅可比”不是麻烦，而是把两类奇异分开看的工具]\label{ins:ak-parallel-twojac}
\cref{eq:ak-parallel-jac} 用 $\mathbf{J}_x,\mathbf{J}_q$ 两个矩阵，恰好把并联的奇异\emph{解耦}成两种独立失效\cite{gosselin1990singularity}：$\mathbf{J}_q$ 退化是一类（驱动侧），$\mathbf{J}_x$ 退化是另一类（末端侧）。串联只有“一个 $\mathbf{J}$ 掉秩”一种奇异；并联因驱动与末端被约束隔开，奇异天然分两侧——这是下一段三类奇异的代数根源。
\end{insight}

\paragraph{三类奇异：并联比串联多一种“危险”。}
基于 $\mathbf{J}_x,\mathbf{J}_q$ 的退化，Gosselin 与 Angeles 给出并联机构奇异的经典三分类\cite{gosselin1990singularity}：

\begin{description}
  \item[第一类：逆运动学奇异（$\det\mathbf{J}_q=0$）] 动平台位于工作空间\emph{边界}（某腿伸到极限、腿与平台共面等），此时末端某速度需\emph{无穷大}驱动速度才能实现——末端\textbf{丢自由度}。这与串联奇异\emph{同性}（末端在某方向瞬时动不了），是“能动范围的边”。
  \item[第二类：直接（正）运动学奇异（$\det\mathbf{J}_x=0$）] 动平台\emph{即便锁死全部驱动}（腿长全固定）仍能\emph{瞬时微动}——平台\textbf{获得一个不受控的自由度}，且此方向上\emph{无法施加/抵抗力}（结构瞬时“撑不住”、刚度塌陷）。\textbf{这是并联独有、串联没有的危险奇异}：发生时平台失控、可能自毁，是并联设计的头号大忌。
  \item[第三类：位形（组合 / 结构）奇异] $\mathbf{J}_x,\mathbf{J}_q$ \emph{同时}退化，源于机构特定几何参数（如某些腿恒共面）的退化构型，须在设计阶段排除。
\end{description}

\begin{pitfall}[陷阱：把并联的“失控奇异”当成串联的“动不了奇异”]
串联运动学奇异的危害是“某方向\emph{推不动}末端”（\cref{def:ak-singularity}，$\det\mathbf{J}=0$）；并联的\emph{第二类}奇异恰恰\textbf{相反}——驱动全锁死，平台\emph{仍会动}，且该方向\emph{撑不住力}。二者方向性截然相反：串联是“失去运动能力”，并联第二类是“失去约束/承载能力”。把并联当串联处理（只防 $\det\mathbf{J}_q=0$、不防 $\det\mathbf{J}_x=0$）会漏掉最致命的失控奇异。这是并联机构安全分析的核心警示：\textbf{逆雅可比奇异管边界，正雅可比奇异管失控，二者都要查}。
\end{pitfall}

\begin{note}[工作空间：并联的“小而满是奇异面”]\label{note:ak-parallel-ws}
并联机构的工作空间显著\emph{小于}同尺度串联臂\cite{siciliano2009robotics,merlet2006parallel}：受各腿行程上下限、被动副（U/S）转角极限、腿间机械干涉三重限制，可达域是这些约束的\emph{交集}，形状复杂且非凸。更棘手的是工作空间\emph{内部}常\emph{穿插着第二类奇异面}，把工作空间割成若干互不连通的\emph{无奇异子域}（singularity-free region）；规划时不仅要保证目标可达，还须保证整条路径\emph{不穿越}第二类奇异面（否则中途失控）。这与串联“工作空间连通、奇异多在边界或孤立面”形成对比，是并联规划（\cref{ch:arm-traj} 的思路可借鉴但约束更严）的特有难点。\textbf{闭链的动力学}（环路闭合约束、Schur 补求解）属另一层面，本书在 \cref{ch:constrained-dynamics} 专门处理；本节只论运动学。
\end{note}

\begin{practice}
\begin{enumerate}
  \item 对 Stewart 平台写出逆运动学 \cref{eq:ak-stewart-ik}，并说明为什么六条腿可\emph{并行独立}求解、而正运动学不能。给出 $\mathbf{J}_q=\mathrm{diag}(\ell_i)$ 的推导（提示：对 $\ell_i^2$ 求导）。
  \item 用一句话各自刻画三类奇异（\cref{thm:ak-parallel-jac} 后）的物理后果，并指出哪一类是串联\emph{没有}的、为什么它最危险。
  \item Delta 机构动平台为何\emph{始终平动}（无姿态变化）？从平行四边形约束解释，并说明这如何把它的正/逆运动学都简化（相比通用 6\nobreakdash-DoF 并联）。
\end{enumerate}
\end{practice}

% ════════════════════════════════════════════════════════════════
\section{机构设计与灵巧度指标}
\label{sec:ak-design}

\paragraph{动机：从“此刻离奇异多远”到“这台臂设计得好不好”。}
\cref{sec:ak-singularity} 的可操作度 $w(\mathbf{q})=\sqrt{\det(\mathbf{J}\mathbf{J}^\top)}$（\cref{eq:ak-manip}）回答的是\emph{给定一台臂、在某个构型}的局部灵活度。但机构\emph{设计}阶段要回答的是另一层问题：\emph{这台臂的几何参数选得好不好？哪种构型最“正”？整个工作空间平均下来够不够灵巧？}本节在可操作度之上，补一组\textbf{设计级}灵巧度指标——条件数 $\kappa(\mathbf{J})$、各向同性、全局条件指数 GCI——它们与 §1.6 的椭球\emph{互补}：椭球量\emph{体积}（能动多快），本节量\emph{形状}（各方向是否均衡）与\emph{全局}（整个工作空间的平均品质）。

\begin{insight}[可操作度椭球量“大小”，条件数量“圆不圆”]\label{ins:ak-design-vs-manip}
两个指标看同一个可操作度椭球（\cref{fig:ak-manip-ellipsoid}）的不同侧面，缺一不可：
\begin{itemize}
  \item \textbf{可操作度} $w=\prod_k\sigma_k$（§1.6）$\propto$ 椭球\emph{体积}：度量“末端整体能动多快”，但\emph{对形状盲}——一个又大又扁的椭球 $w$ 仍可能不小，却在短轴方向几乎动不了。
  \item \textbf{条件数} $\kappa=\sigma_{\max}/\sigma_{\min}$（本节）$=$ 椭球\emph{长短轴之比}：度量“各方向是否\emph{均衡}”，与体积无关。$\kappa=1$ 时椭球为\emph{球}（完美各向同性），$\kappa\to\infty$ 时塌扁（趋奇异）。
\end{itemize}
故 $w$ 大未必“好用”：可能大而偏（某方向极弱）。设计要的是\emph{又大又圆}——$w$ 大\emph{且} $\kappa$ 小。本节正是补上“圆不圆”这一维，与 §1.6 拼成完整的灵巧度图景。
\end{insight}

\begin{definition}[雅可比条件数与各向同性构型]\label{def:ak-condition-number}
雅可比 $\mathbf{J}(\mathbf{q})$ 的\textbf{条件数}定义为最大与最小奇异值之比\cite{angeles2014fundamentals}：
\begin{equation}
  \boxed{\kappa(\mathbf{J})=\frac{\sigma_{\max}(\mathbf{J})}{\sigma_{\min}(\mathbf{J})}\;\ge 1.}
  \label{eq:ak-condition-number}
\end{equation}
$\kappa$ 越接近 $1$ 越好：$\kappa(\mathbf{J})=1$ 当且仅当所有奇异值相等，$\mathbf{J}\mathbf{J}^\top=\sigma^2\mathbf{I}$，可操作度椭球退化为\emph{球}。满足 $\kappa=1$ 的构型称为\textbf{各向同性构型}（isotropic configuration）：末端在\emph{所有方向}上同样灵活、力/速度传递\emph{均衡无偏}，是局部意义下“最正”的位形。$1/\kappa$ 称为\textbf{局部条件指数}（local conditioning index, LCI），取值 $(0,1]$，越大越好。
\end{definition}

\begin{insight}[条件数为何兼管“数值”与“物理”]\label{ins:ak-condition-dual}
$\kappa(\mathbf{J})$ 一身二任\cite{angeles2014fundamentals}：(1) \textbf{数值上}，它正是线性方程 $\dot{\mathbf{x}}=\mathbf{J}\dot{\mathbf{q}}$ 求解（IK 速度级）的\emph{放大因子}——末端速度的相对误差经 $\mathbf{J}^{-1}$ 最多被放大 $\kappa$ 倍，故 $\kappa$ 大则 DiffIK（\cref{eq:ak-dls}）数值病态、对噪声敏感；(2) \textbf{物理上}，它度量\emph{力/速度各向同性}——$\kappa=1$ 时末端在各方向传递力与速度的能力一致，$\kappa$ 大则某方向“能动不能扛、能扛不能动”。设计阶段把臂尽量工作在低 $\kappa$ 区，等于同时优化\emph{数值稳健}与\emph{物理均衡}。这与 §1.6 DLS 的视角衔接：DLS 是\emph{近奇异时}的事后补救，低 $\kappa$ 设计是\emph{从源头}让臂远离病态。
\end{insight}

\begin{theorem}[全局条件指数 GCI]\label{thm:ak-gci}
局部条件数只刻画\emph{单个}构型。要评价一台臂\emph{整个工作空间} $\mathcal{W}$ 的平均品质，Gosselin 与 Angeles 定义\textbf{全局条件指数}（global conditioning index, GCI）为局部条件指数 $1/\kappa$ 在工作空间上的体积平均\cite{angeles2014fundamentals,gosselin1990singularity}：
\begin{equation}
  \boxed{\mathrm{GCI}=\frac{\displaystyle\int_{\mathcal{W}}\frac{1}{\kappa(\mathbf{J}(\mathbf{q}))}\,\mathrm{d}\mathcal{W}}{\displaystyle\int_{\mathcal{W}}\mathrm{d}\mathcal{W}}\;\in(0,1].}
  \label{eq:ak-gci}
\end{equation}
$\mathrm{GCI}$ 越接近 $1$，表示该臂\emph{在整个工作空间普遍接近各向同性}（处处好用）；GCI 小则工作空间内品质悬殊（有些位形很正、有些近奇异）。
\end{theorem}

\begin{insight}[GCI 是“设计目标函数”，可操作度是“运行时指标”]\label{ins:ak-gci-design}
GCI 与可操作度 $w$ 的分工\cite{angeles2014fundamentals}：$w(\mathbf{q})$ 是\emph{运行时}量（随构型实时变化，用于规划避奇异、零空间最大化，见 \cref{sec:ak-singularity} 与 \cref{sec:ak-redundancy}）；GCI 是\emph{设计时}量（对整个工作空间积分、与具体构型无关，是一个\textbf{标量品质数}）。机构尺度综合（dimensional synthesis）正是\emph{以 GCI 为目标函数}，优化连杆长度、关节布置等几何参数，使 GCI 最大——即设计出一台“在它的整个工作空间里都尽量好用”的臂。这把 §1.6 的局部椭球分析\emph{升}到了机构设计层：不再问“此刻好不好”，而问“这套几何参数好不好”。
\end{insight}

\paragraph{量纲陷阱：平移与旋转混在一个雅可比里。}
上述指标有一个\emph{绕不过}的隐患：当任务同时含平移与旋转时，几何雅可比 $\mathbf{J}\in\mathbb{R}^{6\times n}$ 的上三行（角速度，量纲 $\mathrm{rad/s}$，\emph{无量纲}）与下三行（线速度，量纲 $\mathrm{m/s}$，\emph{有长度量纲}）\textbf{量纲不一致}。于是 $\sigma_{\max}/\sigma_{\min}$ 这个比值\emph{依赖于所用单位}（米还是毫米），$\kappa$ 与 GCI 会随单位变化而\emph{失去客观意义}——这是灵巧度度量历史上著名的“量纲不齐”难题\cite{angeles2014fundamentals}。

\begin{theorem}[特征长度归一化]\label{thm:ak-char-length}
解决量纲不齐的标准做法是引入一个\textbf{特征长度}（characteristic length）$L>0$，把雅可比的\emph{平移块除以} $L$（或等价地把姿态以弧度、位置以“$L$ 的倍数”无量纲化），得\emph{齐次化}雅可比\cite{angeles2014fundamentals}：
\begin{equation}
  \mathbf{J}_{\mathrm{hom}}=\begin{bmatrix}\mathbf{J}_\omega\\[2pt] \tfrac{1}{L}\,\mathbf{J}_v\end{bmatrix},
  \label{eq:ak-char-length}
\end{equation}
其中 $\mathbf{J}_\omega,\mathbf{J}_v$ 分别为角速度、线速度子块。$L$ 取使 $\kappa(\mathbf{J}_{\mathrm{hom}})$ \emph{最小化}（或使各向同性可达）的值——此时 $\mathbf{J}_{\mathrm{hom}}$ 的全部奇异值同量纲、$\kappa$ 与 GCI 才\emph{单位无关、客观可比}。
\end{theorem}
\begin{proof}[说明]
$L$ 把线速度块的“米”折算成无量纲数：$\frac{1}{L}\mathbf{J}_v$ 的元素量纲为 $\mathrm{(m/s)/m}=\mathrm{1/s}$，与 $\mathbf{J}_\omega$ 的 $\mathrm{rad/s}=\mathrm{1/s}$ 一致，故 $\mathbf{J}_{\mathrm{hom}}$ 各行同量纲、奇异值可比。几何上 $L$ 给出一个“把转动 $1\,\mathrm{rad}$ 视同平移 $L$ 米”的折算尺度——它把“姿态误差”与“位置误差”放到同一把尺子上。$L$ 的常见取法：使条件数取极小的最优值，或取机构特征尺寸（如可达半径），不同取法对应不同的“位姿距离”度量，须\emph{显式声明}。\textbf{要点}：脱离 $L$ 谈含旋转任务的 $\kappa$/GCI 没有客观意义。
\end{proof}

\begin{pitfall}[陷阱：直接对 $6\times n$ 几何雅可比算条件数]
新手常直接把含平移 + 旋转的 $6\times n$ 几何雅可比丢进 \cref{eq:ak-condition-number} 算 $\kappa$。\textbf{这个数随你用米还是毫米而变}——换单位 $\kappa$ 就变，毫无设计意义。正确流程：(1) 先按 \cref{eq:ak-char-length} 用特征长度 $L$ 把平移块无量纲化；(2) 再算 $\kappa(\mathbf{J}_{\mathrm{hom}})$、GCI；(3) 报告时\emph{声明} $L$ 取值与含义。\emph{纯平移}任务（如 Delta、3\nobreakdash-P 机构）或\emph{纯姿态}任务无此问题（量纲本就齐）；一旦平移、旋转混合，量纲归一是\emph{前提}而非可选。这是灵巧度指标最常被忽视、却最影响结论的一步。
\end{pitfall}

\begin{note}[与 §1.6 的边界：互补不重复]\label{note:ak-design-complement}
本节与 \cref{sec:ak-singularity} 严格\emph{互补}、不重复：§1.6 给\emph{可操作度} $w$（椭球\emph{体积}、运行时避奇异、DLS 事后补救），本节给\emph{条件数族}（椭球\emph{形状}/各向同性、GCI 全局设计度量、量纲归一）。一句话分工：\textbf{$w$ 答“能动多快、离奇异多远”（§1.6），$\kappa$/GCI 答“各方向均不均衡、整台臂设计得正不正”（本节）}。规划/控制（\cref{ch:arm-traj,ch:arm-ctrl}）多用 $w$ 与 DLS；机构尺度综合用 $\kappa$/GCI。读者需要椭球与 DLS 的几何与推导请回看 §1.6，本节不再复述。
\end{note}

\begin{practice}
\begin{enumerate}
  \item 对 2R 平面臂的位置雅可比（\cref{sec:ak-jacobian} 练习），求其各向同性构型（$\kappa=1$）。提示：令 $\mathbf{J}\mathbf{J}^\top\propto\mathbf{I}$，解出 $\theta_2$ 与连杆长 $a_1,a_2$ 应满足的关系（这同时是一个\emph{设计}约束）。
  \item 解释为什么 $w$ 大的构型未必 $\kappa$ 小（举一个“大而扁”椭球的例子），并说明设计上为何要\emph{两者兼顾}。
  \item 一台臂任务含位置（m）+ 姿态（rad）。若不做特征长度归一，把位置单位从 m 改成 mm，$\kappa$ 会如何变化？据此说明 \cref{thm:ak-char-length} 的必要性。
\end{enumerate}
\end{practice}

% ════════════════════════════════════════════════════════════════
\section{工作空间：可达与灵巧}
\label{sec:ak-workspace}

\paragraph{动机：臂“够得到哪些位姿”是规划的前提。}
\cref{sec:ak-ik-analytic} 的逆运动学已多次撞上“目标在工作空间外则无解”（\cref{note:ak-ik-multisol}）。但那里只把“工作空间”当作一个直觉边界一带而过。本节把它\emph{形式化}：一台臂到底能到达哪些末端位姿？这个集合长什么样、边界由什么决定、如何\emph{算}出来？这是任务规划（\cref{ch:arm-traj}）、抓取布置、单元布局的几何前提——目标若不在工作空间内，再好的 IK / 控制器也无能为力。本节把它讲透，并与 \cref{sec:ak-singularity}（奇异 / 可操作度）严格\emph{互补}：那里问“离奇异多远”，这里问“可达域的边在哪、能否任意定姿”。

\paragraph{核心思想：工作空间就是 FK 映射的“像”。}
正运动学（\cref{ch:arm-kin}）是一个映射 $\mathbf{T}_{0e}=f(\mathbf{q})$，把关节空间 $\mathcal{Q}$（受限位约束的关节角集合）送进 $\SE(3)$。臂能到达的全部末端位姿，正是这个映射在\emph{整个可行关节空间}上的\textbf{像集}（image）。于是“工作空间”不是凭空的几何体，而是\emph{由 FK 和关节限位共同决定}的确定对象——这把模糊的“够得到”锚定成可计算的集合。按“是否考虑姿态”，它分三层。

\begin{definition}[三层工作空间]\label{def:ak-workspace}
设 $\mathcal{Q}\subseteq\mathbb{R}^n$ 为满足关节限位的可行关节空间，$f:\mathcal{Q}\to\SE(3)$ 为正运动学映射、$\mathbf{p}(\mathbf{q})\in\mathbb{R}^3$ 为其位置分量\cite{murray1994mathematical}：
\begin{itemize}
  \item \textbf{完整工作空间}（complete workspace）$\mathcal{W}=\{f(\mathbf{q}):\mathbf{q}\in\mathcal{Q}\}\subseteq\SE(3)$：可达的全部\emph{位姿}（含姿态）集合。它是 $\SE(3)$ 的子集，最完备但难以直接可视化。
  \item \textbf{可达工作空间}（reachable workspace）$\mathcal{W}_R=\{\mathbf{p}(\mathbf{q}):\mathbf{q}\in\mathcal{Q}\}\subseteq\mathbb{R}^3$：能以\emph{某一}姿态到达的\emph{位置}集合（忽略姿态）。这是日常说“臂能伸到哪”的那个三维体。
  \item \textbf{灵巧工作空间}（dextrous workspace）$\mathcal{W}_D=\{\mathbf{p}\in\mathbb{R}^3:\forall\mathbf{R}\in\SO(3),\ \exists\,\mathbf{q}\in\mathcal{Q}\ \text{s.t.}\ f(\mathbf{q})=(\mathbf{R},\mathbf{p})\}$：能以\emph{任意}姿态到达的位置集合。
\end{itemize}
三者满足包含关系 $\mathcal{W}_D\subseteq\mathcal{W}_R$，且 $\mathcal{W}_R$ 是 $\mathcal{W}$ 向 $\mathbb{R}^3$ 的投影。
\end{definition}

\begin{insight}[灵巧工作空间只在“机构自由度 $>$ 任务定姿需求”时才非空]\label{ins:ak-ws-dextrous-cond}
灵巧工作空间要求“在该点能摆出\emph{所有} $\SO(3)$ 姿态”，这是个苛刻条件。它\emph{有意义}的前提是机构有\emph{富余}的姿态自由度\cite{murray1994mathematical}：一个 6 自由度臂要在某点实现任意三维姿态，须把 3 个自由度“花”在定姿上、只剩 3 个定位——故 $\mathcal{W}_D$ 一般是 $\mathcal{W}_R$ 的\emph{真子集}（越往可达域边缘，能摆的姿态越少，到边界处往往只剩一种伸直姿态，$\mathcal{W}_D$ 在那里为空）。反过来，\textbf{若任务只需定位（任务自由度 $<$ 机构自由度）}，灵巧工作空间这一概念才真正派上用场——它告诉你“在这些点上，姿态可以随便挑，定位时不必担心够不到某个朝向”。对\emph{平面}臂，“任意姿态”退化为平面内任意转角（$\SO(2)$）；对纯定位任务（无姿态要求），$\mathcal{W}_D$ 的概念不适用，只看 $\mathcal{W}_R$。这与 \cref{sec:ak-redundancy} 的“降维造冗余”同根：自由度相对任务的\emph{富余}，既是零空间的来源，也是灵巧工作空间存在的条件。
\end{insight}

\begin{theorem}[球腕使灵巧域等于可达域]\label{thm:ak-ws-sphwrist}
若臂末端为\textbf{球腕}（三转轴交于一点，\cref{thm:ak-pieper}）且\emph{末端执行器系取在腕心}，则该臂可在任一可达点实现任意姿态，故\cite{murray1994mathematical}
\begin{equation}
  \boxed{\mathcal{W}_D=\mathcal{W}_R,\qquad \mathcal{W}=\mathcal{W}_R\times\SO(3).}
  \label{eq:ak-ws-sphwrist}
\end{equation}
\end{theorem}
\begin{proof}
球腕三轴交于腕心 $\mathbf{p}_w$，其三次转动\emph{不改变}腕心位置（\cref{thm:ak-pieper} 第一步），却可把腕心处的姿态调到任意 $\mathbf{R}\in\SO(3)$（三正交轴张成全部姿态自由度）。于是：只要腕心能到达某点 $\mathbf{p}$（即 $\mathbf{p}\in\mathcal{W}_R$），就能在该点实现任意姿态，故 $\mathbf{p}\in\mathcal{W}_D$；反之 $\mathcal{W}_D\subseteq\mathcal{W}_R$ 恒成立。两向包含即得 $\mathcal{W}_D=\mathcal{W}_R$。此时位置与姿态\emph{解耦}，完整工作空间是可达域与整个 $\SO(3)$ 的笛卡尔积。\emph{前提不可省}：末端系须置于腕心；若有偏置 $d_6\ne0$，末端点绕腕心画球面，边缘处可达姿态受限，$\mathcal{W}_D$ 收缩、上式不再成立（这正是 \cref{sec:ak-ik-analytic} 偏置腕 IK 更难的几何根源）。
\end{proof}

\begin{insight}[又一处“为可解性 / 可用性而设计”]
\cref{thm:ak-ws-sphwrist} 与 \cref{thm:ak-pieper} 的球腕设计动机一脉相承：球腕不仅让 IK 有闭式解（\cref{sec:ak-ik-analytic}），还让\emph{灵巧工作空间最大化到等于可达工作空间}——在可达域内\emph{处处可任意定姿}。这对装配、焊接、码垛等“到点后还要摆特定朝向”的任务极其宝贵：任务规划只需保证目标\emph{位置}在 $\mathcal{W}_R$ 内，姿态自动可行，无需逐点核验姿态可达性。工业 6R 臂几乎清一色球腕，运动学层面的这两条好处（闭式 IK + 灵巧域 $=$ 可达域）是关键推手。
\end{insight}

\paragraph{边界成因：两种本质不同的“边”。}
工作空间的边界（boundary）由两类\emph{机理迥异}的因素造成，区分它们是理解工作空间的关键：

\begin{description}
  \item[可达性边界（关节限位 / 臂展所致）] $\mathcal{W}_R$ 的\emph{外边界}与\emph{内边界}（若有空腔）源于关节走到\emph{行程极限}或连杆\emph{完全伸直 / 完全折叠}。外边界对应臂尽量伸长（$\approx\sum_i a_i$）、内边界对应近端关节折叠形成的不可达空腔。这是“\emph{物理上够不着}”的边——纯粹由几何尺寸与限位决定，与雅可比是否满秩\emph{无必然关系}。
  \item[灵巧性退化（奇异所致）] 在 $\mathcal{W}_R$ \emph{内部}及其边界上，某些点处雅可比\emph{掉秩}（\cref{def:ak-singularity}，\cref{sec:ak-jacobian}）：可达，但末端在某方向瞬时不可动、或无法实现某些姿态。臂\emph{完全伸直}时落在 $\mathcal{W}_R$ 外边界上——这是\emph{边界奇异}（肘奇异，\cref{sec:ak-singularity}），是“可达性边界”与“奇异”\emph{重合}的特例。但\emph{内部}奇异（如腕奇异 $q_5=0$）发生在工作空间深处、远离物理边界，标志的是\emph{灵巧性}退化而非可达性边界——$\mathcal{W}_D$ 的边界正与这类“姿态可达性塌缩”相关。
\end{description}

\begin{pitfall}[陷阱：把“可达性边界”与“奇异”混为一谈]
新手易认为“工作空间边界就是奇异的地方”。\textbf{两者只在外边界（伸直 / 折叠）处重合，其余处各管各的}：(1) 关节限位造成的边界（如某关节转到 $\pm170^\circ$ 停住）处，雅可比可以\emph{满秩、良态}——这是纯粹的可达性边界，无关奇异；(2) 腕奇异 $q_5=0$ 发生在工作空间\emph{内部}，臂在那里\emph{够得到}（不在 $\mathcal{W}_R$ 边界）却\emph{灵巧性退化}（掉一个姿态自由度）。混淆二者会导致错误判断：以为“没到边界就没事”，结果在工作空间深处撞上腕奇异控制飞车（\cref{sec:ak-singularity} 的子雅可比奇异案例）。\textbf{正解}：可达性看\emph{关节限位 + 臂展}（$\mathcal{W}_R$ 的边），灵巧性看\emph{雅可比秩 / 可操作度}（$w\to0$ 的面，\cref{eq:ak-manip}），两套判据并行。
\end{pitfall}

\paragraph{数值计算法：FK 采样 + 边界提取。}
对 2R / 3R 等规整臂，工作空间可\emph{解析}求出（见下例）；但一般臂（异构连杆、复杂限位）的工作空间\emph{无闭式}，工程上用\textbf{蒙特卡洛 / 网格采样}估计——这是最通用、最实用的办法：

\begin{enumerate}
  \item \textbf{采样关节空间}：在可行关节空间 $\mathcal{Q}$ 内（尊重每个 $q_i$ 的限位 $[q_i^{\min},q_i^{\max}]$）均匀随机采样（或网格遍历）大量构型 $\{\mathbf{q}^{(k)}\}_{k=1}^N$。
  \item \textbf{正向映射}：对每个 $\mathbf{q}^{(k)}$ 算 FK（\cref{eq:ak-poschain} 或 POE \cref{eq:ak-poe-space}），得末端位置点云 $\{\mathbf{p}^{(k)}=\mathbf{p}(\mathbf{q}^{(k)})\}$。这堆点近似 $\mathcal{W}_R$（采样越密越准）。
  \item \textbf{边界提取}：对点云求\emph{凸包 / alpha-shape}（非凸时用 alpha-shape 才能捕捉空腔与凹陷），或体素化后取占据体素的外壳，得到 $\mathcal{W}_R$ 的边界面。
  \item \textbf{灵巧域 / 品质叠加}（可选）：若要 $\mathcal{W}_D$，在每个采样点额外核验“能否覆盖一组代表性姿态”；或把可操作度 $w(\mathbf{q}^{(k)})$（\cref{eq:ak-manip}）作为标量\emph{染色}到点云上，得到工作空间的\emph{灵巧度热力图}（哪些区域好用、哪些近奇异）——这把 \cref{sec:ak-singularity} 的局部 $w$ 升为\emph{全工作空间的品质地图}，是单元布局、最优放置（把高频作业区放在高 $w$ 处）的依据。
\end{enumerate}
此法的优点是\emph{对任意臂通用}（只需 FK，不需解析逆解）、易并行（各采样点独立）；缺点是采样估计、边界精度依赖采样密度，且高自由度时 $\mathcal{Q}$ 维数高、需大量样本。

\begin{example}[平面 3R 臂的可达 / 灵巧工作空间——经典 annulus 构造]\label{exam:ak-ws-3r}
取一条平面 3R 臂，连杆长 $a_1>a_2>a_3$，关节角\emph{无限位}（可绕满 $2\pi$）。用\emph{逐层扫掠}（successive sweeping）解析地构造工作空间\cite{murray1994mathematical}（这正是 FK 采样的几何极限版）：
\begin{description}
  \item[第一步] 固定 $\theta_1,\theta_2$，让 $\theta_3$ 扫满 $2\pi$：末端画一个以第二连杆末端为心、半径 $a_3$ 的\emph{圆}。
  \item[第二步] 再让 $\theta_2$ 扫满：上述圆心绕第一连杆末端转，扫出一个\emph{圆环}（annulus），内径 $|a_2-a_3|$、外径 $a_2+a_3$，中心在第一连杆末端。
  \item[第三步] 最后让 $\theta_1$ 扫满：整个圆环绕基座原点转，得\textbf{可达工作空间} $\mathcal{W}_R$——一个以原点为心的圆环，\emph{外径} $a_1+a_2+a_3$（三连杆全伸直）、\emph{内径} $\max(0,\,a_1-a_2-a_3)$（近端折叠形成的空腔；若 $a_1\le a_2+a_3$ 则内径为 $0$、无空腔，$\mathcal{W}_R$ 是实心圆盘）。
\end{description}
\textbf{灵巧工作空间} $\mathcal{W}_D$（能以平面内\emph{任意}转角 $\phi$ 到达的点）更小：要在某点实现任意末端朝向，需第三连杆能绕该点自由转一圈，这要求该点到“前两连杆可达圆环”的距离允许第三连杆任意指向——结果 $\mathcal{W}_D$ 是一个\emph{更窄}的同心圆环，外径 $a_1+a_2-a_3$、内径 $a_1-a_2+a_3$（即相对 $\mathcal{W}_R$ 的外径 $a_1+a_2+a_3$、内径 $a_1-a_2-a_3$，内外边缘各\emph{收窄} $2a_3$）。\textbf{直觉}：最外圈（三连杆伸直）和最内圈处，末端只能以\emph{唯一}朝向到达，无法任意定姿，故被排除出 $\mathcal{W}_D$。这把 $\mathcal{W}_D\subsetneq\mathcal{W}_R$（\cref{def:ak-workspace}）的真包含关系\emph{画了出来}：可达的是大环，可任意定姿的是中间一圈窄环。
\end{example}

\begin{note}[与奇异 / 可操作度（\cref{sec:ak-singularity}）的边界：互补不重复]\label{note:ak-ws-complement}
本节与 \cref{sec:ak-singularity} 严格\emph{互补}：§1.6 给\emph{局部}量——某构型 $\mathbf{q}$ 处雅可比是否掉秩、可操作度 $w(\mathbf{q})$ 离奇异多远；本节给\emph{全局}集合——所有可行 $\mathbf{q}$ 映出的可达 / 灵巧位置域及其边界。二者在\emph{边界奇异}（伸直 / 折叠落在 $\mathcal{W}_R$ 外边界）处衔接：那里既是可达性边界、又是 $w=0$ 的奇异。一句话分工：\textbf{$w$ 答“此刻这台臂在这个构型好不好用”（§1.6），工作空间答“这台臂总共够得到哪些位置、其中哪些能任意定姿”（本节）}。数值上，把 $w(\mathbf{q})$ 染色到工作空间点云，正是用 §1.6 的局部指标给本节的全局集合\emph{上色}——两节由此拼成“局部品质 $+$ 全局范围”的完整图景。读者需要可操作度椭球与 DLS 的几何 / 推导请回看 \cref{sec:ak-singularity}，本节不再复述。
\end{note}

\begin{practice}
\begin{enumerate}
  \item 对 2R 平面臂（连杆长 $a_1,a_2$，无限位），用逐层扫掠（\cref{exam:ak-ws-3r} 的方法去掉第三步）求 $\mathcal{W}_R$：证明它是外径 $a_1+a_2$、内径 $|a_1-a_2|$ 的圆环。何时退化为实心圆盘？2R 臂的 $\mathcal{W}_D$（能以任意平面转角到达的点）是什么？
  \item 给 \cref{exam:ak-ws-3r} 的 3R 臂加关节限位（如 $\theta_1\in[-90^\circ,90^\circ]$）。定性说明 $\mathcal{W}_R$ 如何从整圆环变成\emph{扇环}——这体现了哪一类工作空间边界（\cref{def:ak-workspace} 后的两类成因之一）？
  \item 写出用 FK 采样估计某 6R 臂 $\mathcal{W}_R$ 的伪代码（采样 $\to$ FK $\to$ 取位置 $\to$ 边界），并说明为什么要把可操作度 $w$（\cref{eq:ak-manip}）一并记录、它给规划带来什么信息（提示：\cref{note:ak-ws-complement} 的“灵巧度热力图”）。
  \item 为什么带球腕（末端系在腕心）的臂 $\mathcal{W}_D=\mathcal{W}_R$（\cref{thm:ak-ws-sphwrist}），而 \cref{exam:ak-ws-3r} 的 3R 平面臂 $\mathcal{W}_D\subsetneq\mathcal{W}_R$？二者“定姿能力”的差别在哪？（提示：平面 3R 的“腕”就是第三关节，但末端点\emph{不}在该关节轴上，有 $a_3$ 偏置。）
\end{enumerate}
\end{practice}

% ════════════════════════════════════════════════════════════════
\section{实践小节：库函数与本部实践章指针}
\label{sec:ak-practice}

\paragraph{从公式到代码：Pinocchio 的运动学 API。}
本章的 FK 与雅可比在工程上几乎不手算，而用刚体动力学库（Pinocchio、KDL、Drake）。以 \cref{ch:arm-prac} 所用的 \textbf{Pinocchio} 为例，运动学相关的核心调用映射如下\cite{lynch2017modern}：

\begin{table}[H]\centering\small
\caption{运动学公式 $\to$ Pinocchio API 映射（\cref{ch:arm-prac} 实测使用）}
\label{tab:ak-pinocchio}
\begin{tabular}{p{0.30\textwidth} p{0.32\textwidth} p{0.30\textwidth}}
\toprule
本章概念 & Pinocchio 调用 & 备注 \\
\midrule
正运动学 $\mathbf{T}_{0e}(\mathbf{q})$（\cref{eq:ak-poschain}） & \texttt{forwardKinematics} + \texttt{updateFramePlacements} & 读 \texttt{data.\allowbreak oMf[frame\_id]} 得帧位姿 \\
几何雅可比 $\mathbf{J}$（\cref{eq:ak-jac-columns}） & \texttt{computeFrameJacobian} & 须指定参考系（见下）\\
位置雅可比（取 $\mathbf{J}$ 下 3 行） & \texttt{computeFrameJacobian} 取线速度块 & CBF/避障用（\cref{ch:arm-traj}）\\
\bottomrule
\end{tabular}
\end{table}

\begin{pitfall}[陷阱：雅可比参考系（LOCAL / WORLD / LOCAL\_WORLD\_ALIGNED）]
Pinocchio 的 \texttt{computeFrameJacobian} 按第三个参数 \texttt{reference\_frame} 返回\emph{不同}的雅可比（对应 \cref{note:ak-jac-variants} 的物体/空间/几何变体）：\texttt{LOCAL} 给物体雅可比（twist 在帧自身系）、\texttt{WORLD} 给空间雅可比、\texttt{LOCAL\_WORLD\_ALIGNED} 给“原点在帧、轴平行世界系”的几何雅可比（控制最常用）。\textbf{用错参考系是力控/OSC 最隐蔽的 bug}：$\boldsymbol{\tau}=\mathbf{J}^\top\mathbf{F}$ 里 $\mathbf{J}$ 与 $\mathbf{F}$ 必须在\emph{同一参考系}下，否则力的方向系统性错。务必查清库的默认与你的力/位姿误差表达在哪个系。
\end{pitfall}

\paragraph{承接本部后续。}
本章建立的几何地基贯穿全部：\cref{ch:arm-dyn}（动力学）在雅可比与位姿链上加质量与力、引出操作空间惯量 $\boldsymbol{\Lambda}=(\mathbf{J}\mathbf{M}^{-1}\mathbf{J}^\top)^{-1}$ 与\emph{质量矩阵条件数}（本章 §1.6 已预告其对力矩级 OSC 的影响）；\cref{ch:arm-traj}（轨迹）用本章雅可比做任务空间轨迹与 CBF 反应避障（位置雅可比给 $\nabla h$）；\cref{ch:arm-ctrl}（控制）用 $\boldsymbol{\tau}=\mathbf{J}^\top\mathbf{F}$（\cref{eq:ak-static-duality}）做操作空间控制、阻抗、力位混合，并把 DLS（\cref{eq:ak-dls}）用于 DiffIK；\cref{ch:arm-prac}（实践）把这一切落到 \texttt{arm\_dm} 六轴力矩臂的真工程。完整的工具链与踩坑诊断见 \cref{ch:arm-prac}。

% ════════════════════════════════════════════════════════════════
\section*{常见误解汇总}
\addcontentsline{toc}{section}{常见误解汇总}

\begin{enumerate}
  \item \textbf{“DH 参数表是唯一的。”} 错。标准/改进两套约定、平行轴/相交轴处的自由度、零位选择都使 DH 表不唯一（\cref{note:ak-dh-modified,note:ak-ik-multisol}）。拿到厂商 DH 表必先确认约定与零位。
  \item \textbf{“POE 的零位 $\mathbf{M}$ 随便取一个位姿即可。”} 错。$\mathbf{M}$ 是\emph{所有关节角为零}时的固定末端位姿，旋量 $\mathcal{S}_i$ 也必须在\emph{同一零位}下量（\cref{def:ak-screw} 后的陷阱）。
  \item \textbf{“逆运动学有唯一解。”} 错。6R 带球腕臂一般最多 8 组解（肘/肩/腕组合），还可能无解（工作空间外）。控制时须按最近/避限位/避奇异择一（\cref{note:ak-ik-multisol}）。
  \item \textbf{“雅可比就是一个固定矩阵。”} 错。$\mathbf{J}=\mathbf{J}(\mathbf{q})$ 随构型变，奇异与可操作度都是构型相关的（\cref{def:ak-singularity}）。
  \item \textbf{“近奇异直接用 $\mathbf{J}^{-1}$ / $\mathbf{J}^{+}$ 即可。”} 错。$\mathbf{J}\mathbf{J}^\top$ 病态会让关节速度爆炸，须用 DLS（\cref{eq:ak-dls}）阻尼。但阻尼\emph{在速度级}做最干净，搬到力矩级 OSC 会破坏投影性（\cref{ch:arm-ctrl}）。
  \item \textbf{“表示奇异（万向锁）就是运动学奇异。”} 错。解析雅可比因欧拉角等姿态\emph{表示}退化造成的奇异（\cref{eq:ak-jac-analytic} 的 $\mathbf{T}_\phi$ 退化）是\emph{表示奇异}，可换姿态表示规避；几何雅可比掉秩才是真\emph{运动学奇异}。
  \item \textbf{“冗余只对 $n>6$ 的臂才有。”} 错。即便 6 DoF 臂，只约束部分任务维（如只控位置）就对该任务冗余、有零空间可用（\cref{def:ak-redundancy}）。
\end{enumerate}

% ════════════════════════════════════════════════════════════════
\section*{本章小结}
\addcontentsline{toc}{section}{本章小结}

本章把“关节角 $\leftrightarrow$ 末端位姿”讲透：用位姿链（\cref{eq:ak-poschain}）把 FK 化为齐次变换连乘；两套正运动学建模法——DH 四参数（\cref{eq:ak-dh-matrix}，省参但要放系、分标准/改进）与旋量积 POE（\cref{eq:ak-poe-space}，承 \cref{ch:lie} 指数映射、统一关节类型、只需两系）；逆运动学解析解（Pieper 准则 \cref{thm:ak-pieper}，球腕位置--姿态解耦三步，暴露多解/无解/工作空间；\emph{螺旋式} Paden--Kahan 三子问题 \cref{thm:ak-pk-sub1,thm:ak-pk-sub2,thm:ak-pk-sub3} 把逆解几何化、肘式 6R 四步分解得 8 解 \cref{exam:ak-elbow-ik}；通用偏置腕 6R 用 dialytic 消元、解数尖锐上界 16 \cref{thm:ak-6r-bound}）；雅可比（\cref{eq:ak-jac-columns} 逐列构造、静力学对偶 \cref{eq:ak-static-duality}）；奇异与可操作度（\cref{eq:ak-manip} 椭球、DLS 伪逆 \cref{eq:ak-dls}）；冗余与零空间（resolved-rate \cref{eq:ak-resolved-rate}，复用 \cref{sec:qk-nullspace} 已证的伪逆/投影）。Robot \texttt{arm\_dm} 六轴力矩臂提供了实证血肉：球腕的可解性设计、运动学级零空间 EE 锁 0.00mm 与 $J_1$ 主导自运动、折叠臂子雅可比奇异的诊断。

\begin{table}[H]\centering\small
\caption{本章主要符号}
\label{tab:ak-symbols}
\begin{tabular}{ll}
\toprule
符号 & 含义 \\
\midrule
$\mathbf{q}=(q_1,\dots,q_n)$ & 关节变量（转角 $\theta_i$ / 位移 $d_i$） \\
$\mathbf{T}_{i-1,i}(q_i)$ & 相邻连杆相对位姿（\cref{def:ak-serial}） \\
$\mathbf{T}_{0n}(\mathbf{q})$ & 末端在基座系下位姿（FK，\cref{eq:ak-poschain}） \\
$\theta_i,d_i,a_i,\alpha_i$ & DH 四参数：关节角/偏距/连杆长/扭角（\cref{def:ak-dh}） \\
$\mathcal{S}_i=[\boldsymbol{\omega}_i;\mathbf{v}_i]$ & 关节空间旋量（螺旋轴，\cref{def:ak-screw}） \\
$\mathbf{M}$ & POE 零位末端位姿（\cref{thm:ak-poe-space}） \\
$\mathbf{p}_w$ & 腕心位置（IK 解耦，\cref{thm:ak-pieper}） \\
$\mathbf{J}(\mathbf{q})$ & 几何雅可比（\cref{def:ak-jacobian}） \\
$\mathbf{J}_a,\mathbf{J}_s,\mathbf{J}_b$ & 解析/空间/物体雅可比（\cref{note:ak-jac-variants}） \\
$w(\mathbf{q})=\sqrt{\det(\mathbf{J}\mathbf{J}^\top)}$ & 可操作度（\cref{eq:ak-manip}） \\
$\mathbf{J}^{+},\lambda$ & MP 伪逆、DLS 阻尼系数（\cref{eq:ak-dls}） \\
$\mathbf{N}=\mathbf{I}-\mathbf{J}^{+}\mathbf{J}$ & 零空间投影矩阵（\cref{eq:ak-resolved-rate}） \\
$\boldsymbol{\tau}=\mathbf{J}^\top\mathbf{F}$ & 静力学对偶（\cref{eq:ak-static-duality}） \\
\bottomrule
\end{tabular}
\end{table}

\begin{table}[H]\centering\small
\caption{本章定理与公式速查}
\label{tab:ak-theorems}
\begin{tabular}{lll}
\toprule
名称 & 结论 & 出处 \\
\midrule
位姿链 FK & $\mathbf{T}_{0n}=\prod_i\mathbf{T}_{i-1,i}(q_i)$ & \cref{thm:ak-poschain} \\
标准 DH 矩阵 & 四参数 $\to$ \cref{eq:ak-dh-matrix} & \cref{thm:ak-dh-matrix} \\
空间 POE & $\mathbf{T}=e^{[\mathcal{S}_1]\theta_1}\cdots e^{[\mathcal{S}_n]\theta_n}\mathbf{M}$ & \cref{thm:ak-poe-space} \\
物体 POE & $\mathbf{T}=\mathbf{M}\,e^{[\mathcal{B}_1]\theta_1}\cdots$，$\mathcal{B}_i=\mathrm{Ad}_{\mathbf{M}^{-1}}\mathcal{S}_i$ & \cref{thm:ak-poe-body} \\
Pieper 准则 & 三轴交点/平行 $\Rightarrow$ 闭式 IK & \cref{thm:ak-pieper} \\
Paden--Kahan 子问题 & 绕轴/两圆交/圆球交，$\operatorname{atan2}$ 闭式 & \cref{thm:ak-pk-sub1,thm:ak-pk-sub2,thm:ak-pk-sub3} \\
肘式 6R 解数 & 子问题四步分解 $\Rightarrow$ 至多 8 解 & \cref{exam:ak-elbow-ik} \\
通用 6R 上界 & dialytic 消元 $\Rightarrow$ 至多 16 解（尖锐） & \cref{thm:ak-6r-bound} \\
几何雅可比列 & 转动列 $[\mathbf{z};\mathbf{z}\times(\mathbf{p}_e-\mathbf{p})]$ & \cref{thm:ak-jac-columns} \\
静力学对偶 & $\boldsymbol{\tau}=\mathbf{J}^\top\mathbf{F}$（虚功） & \cref{thm:ak-static-duality} \\
可操作度 & $w=\sqrt{\det(\mathbf{J}\mathbf{J}^\top)}=\prod\sigma_k$ & \cref{def:ak-singularity} \\
DLS 伪逆 & $\dot{\mathbf{q}}=\mathbf{J}^\top(\mathbf{J}\mathbf{J}^\top+\lambda^2\mathbf{I})^{-1}\dot{\mathbf{x}}$ & \cref{thm:ak-dls} \\
resolved-rate & $\dot{\mathbf{q}}=\mathbf{J}^{+}(\dot{\mathbf{x}}_d+\mathbf{K}\mathbf{e})+\mathbf{N}\dot{\mathbf{q}}_{\text{sec}}$ & \cref{thm:ak-resolved-rate} \\
\bottomrule
\end{tabular}
\end{table}

\paragraph{难度标记.}\;
基础（必读）：§1.1 位姿链、§1.2 DH、§1.5 雅可比 $\bigstar$；进阶：§1.3 POE、§1.4 解析 IK、§1.6 奇异 $\bigstar\bigstar$；深入（选读细节）：§1.7 冗余零空间的力矩级对照、偏置腕 IK $\bigstar\bigstar\bigstar$。

% ════════════════════════════════════════════════════════════════
\section*{延伸阅读}
\addcontentsline{toc}{section}{延伸阅读}

\begin{itemize}
  \item \textbf{旋量积 POE 主线}：Lynch \& Park, \emph{Modern Robotics}\cite{lynch2017modern}——本章 §1.3/§1.5 的 POE、空间/物体雅可比、可操作度椭球的权威出处，与本书 \cref{ch:lie} 李群机器一脉相承。
  \item \textbf{DH 与解析逆运动学}：Spong, Hutchinson \& Vidyasagar, \emph{Robot Modeling and Control}\cite{spong2006robot}——§1.2 DH 约定、§1.4 Pieper 准则与球腕解耦三步、§1.6 奇异分类的经典教材。
  \item \textbf{李群几何基础}：本书 \cref{ch:lie}（指数映射 \cref{thm:se3-exp}、伴随 \cref{eq:Ad}、位姿运动学 \cref{thm:se3-kinematics}）是 POE 与雅可比的数学地基。
  \item \textbf{零空间/伪逆/DLS 的完整证明}：本书 \cref{sec:qk-nullspace}（四足部）用拉格朗日乘子法证 MP 伪逆最小范数、零空间投影、阻尼最小二乘——本章复用其结论。
  \item \textbf{螺旋式逆解与消元法}：Murray--Li--Sastry, \emph{A Mathematical Introduction to Robotic Manipulation}\cite{murray1994mathematical}——\cref{sec:ak-ik-paden-kahan} 的 Paden--Kahan 子问题、\cref{sec:ak-ik-elimination} 的 dialytic 消元与 16 解上界的权威出处；原始消元工作见 Raghavan--Roth\cite{raghavan1993inverse}、Manocha--Canny（实时化为广义特征值问题）\cite{manocha1994realtime}、Lee--Liang（16 尖锐界）\cite{lee1988displacement}。
  \item \textbf{偏置腕/统一 IK}：近期工作用子问题分解 / 消元法处理非球腕臂的半解析 IK（IK-Geo 等），其经典代数地基即 \cref{sec:ak-ik-elimination} 的 dialytic 消元。
\end{itemize}

% ════════════════════════════════════════════════════════════════
\section*{本章与后续的关系}
\addcontentsline{toc}{section}{本章与后续的关系}

\begin{itemize}
  \item \textbf{$\to$ \cref{ch:arm-dyn}（动力学）}：在位姿链与雅可比上加质量/惯量，引出 $\mathbf{M}(\mathbf{q})\ddot{\mathbf{q}}+\mathbf{C}\dot{\mathbf{q}}+\mathbf{g}=\boldsymbol{\tau}$、操作空间惯量 $\boldsymbol{\Lambda}=(\mathbf{J}\mathbf{M}^{-1}\mathbf{J}^\top)^{-1}$、以及本章 §1.6 预告的质量矩阵条件数。
  \item \textbf{$\to$ \cref{ch:arm-traj}（轨迹与避障）}：用本章雅可比做任务空间轨迹、用位置雅可比给 CBF 反应避障的 $\nabla h$；冗余零空间亦可用于避障。
  \item \textbf{$\to$ \cref{ch:arm-ctrl}（控制）}：用静力学对偶 $\boldsymbol{\tau}=\mathbf{J}^\top\mathbf{F}$ 做操作空间控制、阻抗、力位混合；DLS 用于 DiffIK；§1.6 的奇异/条件数告诫贯穿其中。
  \item \textbf{$\to$ \cref{ch:arm-prac}（实践）}：本章 FK/Jacobian 经 Pinocchio API 落到 \texttt{arm\_dm} 六轴力矩臂全栈。
  \item \textbf{$\to$ \cref{ch:rlmc-manip}（P8 强化学习操作）}：本章为其补经典地基——FK/Jacobian（\cref{sec:ak-jacobian}）、DiffIK 的 DLS 伪逆背景（\cref{sec:ak-singularity}）此前在 P8 仅给公式不给来历，本章补全。
  \item \textbf{对照 \cref{ch:quad_kin}（P7 四足）}：本章给定基串联臂通用框架，P7 给规整短链解耦捷径；零空间理论 \cref{sec:qk-nullspace} 共享。
\end{itemize}

% ════════════════════════════════════════════════════════════════
\section*{故障排查}
\addcontentsline{toc}{section}{故障排查}

\begin{table}[H]\centering\small
\caption{机械臂运动学常见故障 $\to$ 根因 $\to$ 对策}
\label{tab:ak-troubleshoot}
\begin{tabularx}{\linewidth}{p{0.27\textwidth} L L}
\toprule
现象 & 可能根因 & 对策 \\
\midrule
FK 算出的末端位姿与实物系统性偏差 & DH 约定混用（标准 vs 改进）、零位定义错、URDF 与代码模型不一致 & 确认 DH 约定与零位（\cref{note:ak-dh-modified}）；以 URDF 为准重建 POE 旋量 \\
POE 的 FK 全错 & 零位 $\mathbf{M}$ 或旋量 $\mathcal{S}_i$ 未在同一零位下量 & 重取统一零位（\cref{def:ak-screw} 后陷阱） \\
IK 求解器“无解” & 目标在工作空间外、姿态不可达 & 检查目标在可达域内（\cref{note:ak-ik-multisol}）；放宽姿态约束 \\
IK 选到“怪异”构型（突然翻腕/换肘） & 多解中跳到了非最近解 & 按关节空间最近解 + 避限位/避奇异择解 \\
近某构型关节速度暴涨/控制飞车 & 运动学奇异，$\mathbf{J}\mathbf{J}^\top$ 病态 & 用 DLS（\cref{eq:ak-dls}）加阻尼；规划避奇异（最大化 $w$） \\
力控/OSC 力方向系统性错 & 雅可比参考系用错（LOCAL/WORLD/对齐） & 确保 $\mathbf{J}$ 与 $\mathbf{F}$ 同参考系（\cref{sec:ak-practice} 陷阱） \\
按子雅可比切分做力位控制时飞车 & 折叠构型下子雅可比奇异 & 用完整雅可比 + 冗余进零空间（\cref{ch:arm-ctrl}） \\
力矩级零空间 EE 漂移大 & 质量矩阵病态，正则破坏投影性 & 冗余解析改用运动学/速度级（\cref{thm:ak-resolved-rate}）；条件数分析见 \cref{ch:arm-dyn} \\
\bottomrule
\end{tabularx}
\end{table}
